\chapter{Gry nieskończone i aksjomat determinacji}\label{ch:determinacy}
\chaptermark{Gry nieskończone i determinacja}
\section{Gra o doskonałej informacji}
W skończonej grze można często wyznaczyć zwycięzcę przez indukcję wsteczną. Gry o nieskończenie wielu ruchach pokazują, jak bardzo zmienia się sytuacja po przejściu do ciągów nieskończonych.
\begin{definicja}[Gra $G(A)$]\label{def:game-a}
Niech $A\subseteq\omega^\omega$, gdzie $\omega^\omega$ jest zbiorem wszystkich funkcji z $\omega$ w $\omega$. Dwaj gracze, I i II, na przemian wybierają liczby naturalne $a_0,a_1,a_2,\ldots$; gracz I wykonuje ruchy o indeksach parzystych, a II o nieparzystych. Gracz I wygrywa dokładnie wtedy, gdy otrzymany ciąg $a\in\omega^\omega$ należy do $A$. W przeciwnym razie wygrywa II.
\end{definicja}
\begin{definicja}[Strategia i determinacja]\label{def:strategy-ad}
Strategią gracza I jest funkcja przypisująca każdemu skończonemu ciągowi ruchów o parzystej długości jego następny ruch. Strategią gracza II jest funkcja przypisująca każdemu skończonemu ciągowi o nieparzystej długości jego następny ruch. Strategia jest wygrywająca, jeżeli zapewnia zwycięstwo przeciw każdemu możliwemu ciągowi odpowiedzi przeciwnika. Gra jest \emph{zdeterminowana}, jeśli któryś z graczy ma strategię wygrywającą. Aksjomat determinacji $\mathrm{AD}$ głosi, że każda gra $G(A)$ dla $A\subseteq\omega^\omega$ jest zdeterminowana.
\end{definicja}
\begin{przyklad}[Wymuszenie pierwszego ruchu]
Niech $A=\{a\in\omega^\omega:a(0)=0\}$. Gracz I wygrywa, wybierając w pierwszym ruchu zero; dalsze ruchy nie mają znaczenia. Jeśli $A=\{a:a(1)=0\}$, analogicznie wygrywa gracz II.
\end{przyklad}
\section{Gry otwarte i zamknięte}
\begin{definicja}[Podstawowe otoczenia w przestrzeni Baire'a]\label{def:baire-cylinders}
Dla skończonego ciągu $s\in\omega^{<\omega}$ definiujemy cylinder $[s]=\{a\in\omega^\omega:s\subseteq a\}$. Zbiór $U\subseteq\omega^\omega$ jest otwarty, jeżeli dla każdego $a\in U$ istnieje skończony początkowy fragment $s\subseteq a$ taki, że $[s]\subseteq U$. Zbiór jest zamknięty, jeżeli jego dopełnienie jest otwarte.
\end{definicja}
\begin{twierdzenie}[Determinacja gier otwartych; twierdzenie Gale'a--Stewarta]\label{thm:open-determinacy}
Już w ZF, bez aksjomatu wyboru, jeśli
$A\subseteq\omega^\omega$ jest otwarty,
to gra $G(A)$ jest zdeterminowana.
\end{twierdzenie}
\begin{proof}
Niech $T:=\omega^{<\omega}$ będzie zbiorem
wszystkich skończonych pozycji gry.
Gracz I ma ruch po pozycjach parzystej
długości, gracz II po nieparzystej.
Dla $s\in T$ piszemy
$s^\frown n$ na oznaczenie pozycji
przedłużonej o jeden ruch $n$.

\emph{Krok 1: rekurencyjny zbiór pozycji wygrywanych.}
Niech
$W_0:=\{s\in T:[s]\subseteq A\}$.
Dla dowolnego $W_\alpha\subseteq T$
definiujemy
\[
\begin{split}
W_{\alpha+1}:=W_\alpha
 &\cup\{s\in T:|s|\text{ parzyste},
       \ \exists n\in\omega\
                    (s^\frown n\in W_\alpha)\}\\
 &\cup\{s\in T:|s|\text{ nieparzyste},
       \ \forall n\in\omega\
                    (s^\frown n\in W_\alpha)\}.
\end{split}
\]
Dla niezerowych ordynałów granicznych
$\lambda$ kładziemy
$W_\lambda:=\bigcup_{\alpha<\lambda}W_\alpha$.
Wszystkie wyrazy są podzbiorami $T$,
a operacja przejścia jest jednoznaczna.
Dlatego Twierdzenie~\ref{thm:ordinal-recursion}
zapewnia istnienie rodziny na dowolnym
ustalonym odcinku ordynałów, bez AC.

\emph{Krok 2: stabilizacja.}
Z Twierdzenia Hartogsa~\ref{thm:choice-hartogs}
dla zbioru $\mathcal P(T)$ mamy
dobrze uporządkowany zbiór, którego nie
można różnowartościowo odwzorować
w $\mathcal P(T)$. Biorąc jego
typ porządkowy (konstrukcja w dowodzie
Twierdzenia~\ref{thm:initial-exists},
niewymagająca wyboru dla już
dobrze uporządkowanego zbioru),
otrzymujemy ordynał $\theta$
bez iniekcji $\theta\to\mathcal P(T)$.
Skonstruujmy $W_\alpha$ dla
$\alpha\le\theta$.
Gdyby $W_\alpha\ne W_{\alpha+1}$
dla wszystkich $\alpha<\theta$,
monotoniczność rodziny dawałaby
iniekcję
$\alpha\mapsto W_\alpha$
ze zbioru $\theta$ do $\mathcal P(T)$.
To sprzeczność.
Istnieje więc $\eta<\theta$
z $W_\eta=W_{\eta+1}$,
a rekurencja daje tę samą wartość
na wszystkich dalszych etapach.
Niech $W:=W_\eta$.
Dla każdego $s\in W$ niech
$\rho(s)$ będzie najmniejszym
$\alpha\le\eta$ takim, że
$s\in W_\alpha$.
Funkcja rang istnieje przez
jednoznaczność i zastępowanie.

\emph{Krok 3: jeśli $\varnothing\in W$,
wygrywa I.}
Jeśli $s\in W_0$, to
$[s]\subseteq A$, więc każdy
wynik z przedłużeniem $s$ jest w $A$.
Dla $s\in W\setminus W_0$
ranga $\rho(s)$ jest następnikiem
$\alpha+1$, bo na etapach granicznych
nie pojawiają się nowe elementy.
Jeśli ruch ma gracz I,
istnieje $n\in\omega$
z $s^\frown n\in W_\alpha$.
I wybiera \emph{najmniejszego}
takiego $n$.
Jeżeli ruch ma II, z definicji
$W_{\alpha+1}$ każdy następnik
$s^\frown n$ ma rangę najwyżej $\alpha$.
Na wszystkich ruchach, dopóki
nie osiągnięto $W_0$, ranga ściśle maleje.
Nie istnieje nieskończony ściśle malejący
ciąg ordynałów: zbiór jego wartości
miałby element najmniejszy,
a następny wyraz byłby mniejszy.
Zatem po skończenie wielu ruchach
osiągamy $W_0$ i I wygrywa.
Dalsze ruchy I można ustalić jako $0$.

\emph{Krok 4: jeśli $\varnothing\notin W$,
wygrywa II.}
Dla pozycji $s\notin W$, z której
ruch ma I, każdy jej następnik
leży poza $W$: w przeciwnym razie
jego pewna ranga wprowadziłaby
$s$ do $W$.
Jeżeli ruch ma II, istnieje
przynajmniej jeden następnik
$s^\frown n\notin W$.
Gdyby wszystkie następniki należały
do $W$, ich rangi utworzyłyby
zbiór ordynałów. Kładąc
\[
\beta:=\sup_{n\in\omega}
              (\rho(s^\frown n)+1),
\]
otrzymalibyśmy
$s^\frown n\in W_\beta$
dla wszystkich $n$, a więc
$s\in W_{\beta+1}\subseteq W$,
sprzeczność (stabilność rodziny).
II wybiera najmniejszego $n$
z $s^\frown n\notin W$.
Cała rozgrywka pozostaje poza $W_0$.
Gdyby jej wynik należał do
otwartego $A$, pewien jego
skończony prefiks $s$ spełniałby
$[s]\subseteq A$,
a więc $s\in W_0$.
To niemożliwe, zatem II wygrywa.

Obie strategie są funkcjami
zadanymi jednoznacznie przez
najmniejszy dopuszczalny ruch
w $\omega$; w żadnym miejscu
nie zastosowano AC.
\end{proof}

\begin{przyklad}[Dwa kroki i kolejność kwantyfikatorów]
\label{ex:ad-two-step-quantifiers}
Niech
$A=\{a\in\omega^\omega:a(0)=0
\text{ oraz }a(1)\ne0\}$.
Jeżeli I w pierwszym ruchu wybierze
$a(0)\ne0$, przegrywa niezależnie od
odpowiedzi przeciwnika. Gdy wybierze
$a(0)=0$, gracz II odpowiada $a(1)=0$
i również zapewnia wypłatę spoza $A$.
Dlatego II ma strategię wygrywającą:
w pierwszym swoim ruchu gra $0$,
a później także może grać $0$.

Warunek dwóch pierwszych ruchów ma postać
\[
\forall u\in\omega\;\exists v\in\omega\;
   (u\ne0\;\lor\;v=0).
\]
Nie wystarcza więc, że istnieje para
ruchów dająca wypłatę należącą do $A$:
gracz I musi mieć odpowiedź zwycięską
na \emph{każdy} możliwy ruch przeciwnika.
\end{przyklad}
\begin{figure}[htbp]
\centering
\begin{tikzpicture}[
  level distance=13mm,
  every node/.style={font=\small,align=center},
  edge from parent/.style={draw,-latex},
  level 1/.style={sibling distance=70mm},
  level 2/.style={sibling distance=30mm}
]
\node {$\varnothing$}
  child {node {$a(0)=0$}
    child {node {$a(1)=0$\\II wygrywa}
      edge from parent node[left] {\scriptsize II}}
    edge from parent node[left] {\scriptsize I}}
  child {node {$a(0)\ne0$\\II wygrywa}
    edge from parent node[right] {\scriptsize I}};
\end{tikzpicture}
\caption{Drzewko dwóch ruchów: na każdą decyzję
gracza I istnieje odpowiedź zapewniająca
zwycięstwo graczowi II. Późniejsze ruchy
nie zmieniają wyniku.}
\label{fig:ad-two-moves}
\end{figure}

\begin{uwaga}[Dlaczego AD nie jest aksjomatem ZFC?]\label{rem:ad-ac-incompat}
Pełny aksjomat determinacji jest sprzeczny z AC: w ZFC istnieją zbiory wypłat, dla których żaden gracz nie ma strategii wygrywającej. Ścisły dowód wymaga konstrukcji z wykorzystaniem dobrego uporządkowania przestrzeni strategii; przedstawiamy ją poniżej. AD rozważa się zatem przede wszystkim jako dodatkowy aksjomat nad ZF, nie nad ZFC.
\end{uwaga}
\begin{twierdzenie}[AC implikuje istnienie gry niezdeterminowanej]\label{thm:ac-not-ad}
W ZFC istnieje zbiór $A\subseteq\omega^\omega$ taki, że gra $G(A)$ nie jest zdeterminowana.
\end{twierdzenie}
\begin{proof}
Niech $\mathfrak c=|\omega^\omega|$. Z AC dobrze porządkujemy wszystkie strategie obu graczy, łącznie jako rodzinę długości $\mathfrak c$. Dla każdej strategii gracza I wybierzemy rozgrywkę zgodną z tą strategią, która zakończy się przegraną I; dla każdej strategii gracza II wybierzemy rozgrywkę zgodną z nią, która zakończy się wygraną I. Wykonujemy te wybory rekurencyjnie po początkowym ordynale mocy $\mathfrak c$, zachowując dwie rozłączne rodziny rozgrywek $L$ i $H$.

Na etapie $\alpha<\mathfrak c$ ustalona strategia ma $\mathfrak c$ zgodnych z nią rozgrywek: przeciwnik może dowolnie wybierać nieskończony ciąg swoich ruchów. Wcześniej oznaczono mniej niż $\mathfrak c$ rozgrywek (dla nieskończonej liczby kardynalnej $\mathfrak c$ i $\alpha<\mathfrak c$ zachodzi $|\alpha|\cdot\aleph_0<\mathfrak c$), więc można wybrać świeżą rozgrywkę. Dla strategii I umieszczamy ją w $L$, dla strategii II w $H$. Wybór świeżego elementu ustalamy przez globalny dobry porządek $\omega^\omega$. Na końcu niech $A=H$. Każda strategia I ma zgodną rozgrywkę w $L\subseteq\omega^\omega\setminus A$, więc nie wygrywa zawsze. Każda strategia II ma zgodną rozgrywkę w $H=A$, więc również nie wygrywa zawsze. Gra $G(A)$ jest niezdeterminowana.
\end{proof}
\begin{uwaga}[Granice przedstawionej metody]
Determinacja gier otwartych nie oznacza determinacji wszystkich gier. Dowód powyżej wykorzystuje pełny AC; nie wolno go przenosić do ZF+AD. W kolejnych rozdziałach omówimy konstrukcję modeli wewnętrznych i wymuszanie jako metody precyzyjnego badania niezależności aksjomatów.
\end{uwaga}

\section{Przykład gry zamkniętej}
\begin{przyklad}[Warunek bezpieczeństwa]
Niech $A=\{a\in\omega^\omega:\forall n\in\omega\ a(2n)=0\}$. Zbiór $A$ jest zamknięty: jeśli $a\notin A$, to pewien parzysty wyraz $a(2n)$ jest niezerowy, a każdy ciąg o tym samym prefiksie długości $2n+1$ również leży poza $A$. Gracz I wygrywa, zawsze grając zero. Ten przykład pokazuje, że zamkniętość warunku wypłaty nie przesądza o tym, który gracz wygrywa.
\end{przyklad}
\begin{wniosek}[Determinacja gier zamkniętych]\label{cor:closed-determinacy}
W ZF każda gra $G(A)$ z zamkniętym
$A\subseteq\omega^\omega$ jest zdeterminowana.
\end{wniosek}
\begin{proof}
Niech $B:=\omega^\omega\setminus A$,
który jest otwarty.
Konstruujemy zbiory $V_\alpha$
tak samo jak w dowodzie
Twierdzenia~\ref{thm:open-determinacy},
lecz ze \emph{zamienionymi rolami}:
\[
V_0:=\{s:[s]\subseteq B\}.
\]
Na pozycji, w której ruch ma II,
włączamy $s$ do $V_{\alpha+1}$
po istnieniu choć jednego następnika
w $V_\alpha$;
na ruchu I wymagamy, aby
\emph{wszyscy} następnicy
należeli do $V_\alpha$.
Zachowujemy już uzyskane pozycje,
a na etapach granicznych bierzemy
ich sumy.
Z Twierdzenia Hartogsa i rekurencji
wynika stabilizacja dokładnie
jak dla gier otwartych.

Jeśli pozycja pusta należy
do stabilnego zbioru, II
na swoich ruchach wybiera
najmniejszego następnika
o mniejszej randze.
Na ruchach I wszystkie następniki
mają rangę mniejszą.
Dobrze ufundowane rangi wymuszają
osiągnięcie pozycji $s\in V_0$,
a wtedy każdy wynik gry leży w $B$
i wygrywa II.

W przeciwnym przypadku I
na swoich ruchach wybiera
najmniejszego następnika
poza stabilnym zbiorem,
a na ruchach II wszystkie
następniki pozostają poza nim.
Wynik nie może należeć do
otwartego $B$, gdyż miałby
prefiks w $V_0$.
Zatem wynik należy do $A$
i wygrywa I.
Wszystkie wybory są kanoniczne
i wykonane w ZF.
\end{proof}


\section{Własność Baire'a, determinacja projektywna i ich następstwa}
\label{sec:ad-regularity}

\subsection{Z jakich założeń wolno korzystać?}
W rozdziale~\ref{ch:determinacy} zdefiniowaliśmy grę $G(A)$,
strategię wygrywającą i aksjomat determinacji $\mathrm{AD}$.
Dotychczas pracowaliśmy zwykle w ZFC. Teraz wyraźnie zmieniamy
teorię bazową: przy rozważaniu pełnego $\mathrm{AD}$ pracujemy
w $\mathrm{ZF}+\mathrm{AD}$, \emph{bez} aksjomatu wyboru.
Twierdzenie~\ref{thm:ac-not-ad} pokazało, że $\mathrm{AD}$ i
$\mathrm{AC}$ nie mogą obowiązywać jednocześnie.

Nie wolno utożsamiać trzech zdań:
\begin{enumerate}
\item wszystkie gry $G(A)$ na liczbach naturalnych są zdeterminowane ($\mathrm{AD}$);
\item gry $G(A)$ są zdeterminowane tylko dla \emph{projektywnych} zbiorów wypłat ($\mathrm{PD}$);
\item każda gra z wypłatą otwartą jest zdeterminowana
      (twierdzenie Gale'a--Stewarta, udowodnione wcześniej).
\end{enumerate}
Pierwsze zdanie jest niezgodne z $\mathrm{AC}$; drugie może być
rozważane wraz z ZFC i wynika z odpowiednio silnych założeń
o dużych liczbach kardynalnych; trzecie jest zwykłym twierdzeniem ZF
(Twierdzenie~\ref{thm:open-determinacy}
i Wniosek~\ref{cor:closed-determinacy},
wykazane bez AC).

\begin{lemat}[AD pozwala wybierać z przeliczalnej rodziny zbiorów rzeczywistych]
\label{lem:ad-countable-choice-reals}
W $\mathrm{ZF}+\mathrm{AD}$ dla każdej funkcji
$n\mapsto X_n$ ($n\in\N$), której wartości są
niepustymi podzbiorami $\N^\N$, istnieje funkcja
$n\mapsto x_n$ taka, że $x_n\in X_n$ dla każdego $n$.
To jest \emph{wybór przeliczalny dla zbiorów liczb rzeczywistych},
a nie pełny aksjomat wyboru.
\end{lemat}
\begin{proof}
Zdefiniujmy grę o ruchach będących liczbami naturalnymi.
W swoim pierwszym ruchu gracz I wybiera indeks $n$.
Pozostałe ruchy gracza I nie mają znaczenia dla wypłaty.
Kolejne ruchy gracza II tworzą nieskończony ciąg $x$.
Gracz II wygrywa wtedy i tylko wtedy, gdy $x\in X_n$.

Nie istnieje wygrywająca strategia gracza I.
Istotnie, każda jego ustalona strategia na początku
wyznacza konkretny indeks $n$. Z niepustości $X_n$
możemy wybrać pojedynczy ciąg $x\in X_n$ i polecić
II rozegrać go współrzędna po współrzędnej, niezależnie
od dalszych ruchów I. Ta odpowiedź pokonuje rozważaną
strategię I. Wybór jednego świadka po ustaleniu $n$
nie jest zastosowaniem nieskończonej zasady wyboru.
Z AD wynika, że II ma strategię wygrywającą $\tau$.
Dla każdego $n$ ustalmy rozgrywkę, w której I gra
najpierw $n$, a potem same zera. Kolejne odpowiedzi
$\tau$ tworzą ciąg $x_n\in X_n$. Przyporządkowanie
$n\mapsto x_n$ jest funkcją wyboru; zdefiniowaliśmy
ją jednym wzorem z użyciem tej samej strategii $\tau$.
\end{proof}

\subsection{Zbiory pierwszej kategorii i własność Baire'a}
\begin{definicja}[Zbiory nigdziegęste i pierwszej kategorii]
\label{def:ad-meager}
Niech $X$ będzie przestrzenią topologiczną.
Zbiór $F\subseteq X$ jest \emph{nigdziegęsty}, jeżeli
wnętrze jego domknięcia jest puste:
$\operatorname{int}_X\overline F=\varnothing$.
Zbiór $M\subseteq X$ jest \emph{pierwszej kategorii},
inaczej \emph{mizerny}, jeśli istnieją zbiory nigdziegęste
$F_n\subseteq X$ dla $n\in\N$ takie, że
$M\subseteq\bigcup_{n\in\N}F_n$.
Zbiór jest \emph{rezydualny},
jeżeli jego dopełnienie jest pierwszej kategorii.
\end{definicja}
Przez \emph{rezydualny} rozumiemy więc ``dopełnienie małe
w sensie kategorii'', a nie ``dopełnienie jest przeliczalne''.
Dla przykładu $\Q$ jest pierwszej kategorii w $\R$,
gdyż jest przeliczalną sumą singletonów, z których każdy
jest nigdziegęsty; $\R\setminus\Q$ jest rezydualny.

\begin{definicja}[Własność Baire'a]
\label{def:ad-baire-property}
Zbiór $A\subseteq X$ ma \emph{własność Baire'a}, gdy
istnieje zbiór otwarty $U\subseteq X$ taki, że
$A\mathbin{\triangle}U=(A\setminus U)\cup(U\setminus A)$
jest pierwszej kategorii w $X$.
\end{definicja}
W teorii ZF przeliczalna suma zbiorów pierwszej kategorii
nie musi automatycznie być pierwszej kategorii: może
zabraknąć jednoczesnego wyboru ich przeliczalnych pokryć.
W naszych dowodach pod AD uzasadnia ją
Lemat~\ref{lem:ad-meager-countable-union}:
wykorzystujemy wcześniej dowiedziony
wybór przeliczalny dla zbiorów
liczb rzeczywistych, a nie pełny AC.


\begin{lemat}[Przeliczalna suma zbiorów pierwszej kategorii w AD]
\label{lem:ad-meager-countable-union}
W teorii ZF+$\mathrm{AD}$, jeśli
$(M_n)_{n\in\omega}$ jest ciągiem
podzbiorów przestrzeni Baire'a
$\omega^\omega$ pierwszej kategorii,
to $\bigcup_{n\in\omega}M_n$
jest pierwszej kategorii.
Nie zakładamy pełnego AC ani DC.
\end{lemat}
\begin{proof}
\emph{Krok 1: kody domkniętych zbiorów.}
Niech $T:=\omega^{<\omega}$.
Ponieważ $T$ ma ustaloną enumerację
(Twierdzenie~\ref{thm:card-n2}),
każdy jego podzbiór można jednoznacznie
zakodować elementem $2^\omega$.
Dla domkniętego $F\subseteq\omega^\omega$
połóżmy
\[
T_F:=\{s\in T:[s]\cap F\ne\varnothing\}.
\]
Jest to zbiór istniejący przez separację.
Dla $x\in\omega^\omega$ mamy
\[
x\in F\quad\Longleftrightarrow\quad
\forall k\in\omega\;(x|_k\in T_F).
\]
Kierunek wprost jest natychmiastowy;
w przeciwnym razie otwartość
$\omega^\omega\setminus F$ daje
cylinder $[x|_k]$ rozłączny z $F$.
Zatem kod $T_F$ jednoznacznie wyznacza $F$.
Ponadto $F$ jest nigdziegęsty dokładnie
wtedy, gdy żaden cylinder nie zawiera się
w $F$, co przy domknięciu $F$ daje
\[
\forall s\in T\;\exists t\sqsupseteq s
\quad t\notin T_F.
\]
Rzeczywiście z pustego wnętrza
domknięcia otrzymujemy dla każdego $[s]$
mniejszy cylinder $[t]$ rozłączny z $F$.

\emph{Krok 2: kod świadectwa małości.}
Dla każdego $M_n$ istnieje ciąg
$(F_{n,j})_{j\in\omega}$ zbiorów
nigdziegęstych pokrywających $M_n$.
Każdy z nich można zastąpić
domknięciem, które nadal jest nigdziegęste.
Zakodujmy ciąg kodów ich drzew
jednym elementem $2^\omega$,
korzystając ze zdefiniowanej wcześniej
bijekcji $\omega\times\omega\to\omega$
(Twierdzenie~\ref{thm:card-n2}).
Niech $C_n\subseteq2^\omega$ będzie
zbiorem wszystkich poprawnych kodów
takich świadectw pokrycia $M_n$.
Jest on niepusty, a przez separację
w $2^\omega$ istnieje w ZF.

\emph{Krok 3: zastosowanie AD.}
Lemat~\ref{lem:ad-countable-choice-reals}
daje jedną funkcję $n\mapsto c_n$
z $c_n\in C_n$ dla wszystkich $n$.
Dekodując $c_n$ otrzymujemy podwójnie
indeksowaną rodzinę
$(F_{n,j})_{(n,j)\in\omega^2}$
domkniętych zbiorów nigdziegęstych
z
\[
\bigcup_{n\in\omega}M_n
 \subseteq\bigcup_{(n,j)\in\omega^2}F_{n,j}.
\]
Prawa strona jest przeliczalną sumą
zbiorów nigdziegęstych, ponieważ
$\omega^2$ jest jawnie przeliczalny.
Zatem $\bigcup_nM_n$ jest pierwszej kategorii.
Nie użyto AC poza konsekwencją AD
z wcześniejszego lematu.
\end{proof}

Nie należy mylić tej własności \emph{zbioru} z twierdzeniem
Baire'a o \emph{przestrzeni}: w przestrzeni zupełnej
metrycznie żaden niepusty zbiór otwarty nie jest
pierwszej kategorii w tej przestrzeni.

\begin{lemat}[Postać konstrukcyjna twierdzenia Baire'a]
\label{lem:ad-baire-category}
Niech $X=\N^\N$ z topologią, której bazą są cylindry
$[s]=\{x\in\N^\N:s\text{ jest początkowym fragmentem }x\}$,
$s\in\N^{<\omega}$. Jeżeli $D_n$ ($n\in\N$) są otwarte
i gęste w $X$, to $\bigcap_nD_n$ przecina każdy niepusty
cylinder $[s]$. W szczególności żaden cylinder $[s]$
nie jest pierwszej kategorii.
\end{lemat}
\begin{proof}
Ustalmy $[s]$. Wybieramy najpierw najmniejsze w ustalonej
enumeracji skończonych ciągów przedłużenie $s_0\supseteq s$
takie, że $[s_0]\subseteq D_0$; istnieje ono, bo $D_0$
jest otwarty i gęsty. Następnie indukcyjnie wybieramy
najmniejsze przedłużenie $s_{n+1}\supseteq s_n$ długości
większej niż $n+1$, dla którego $[s_{n+1}]\subseteq D_{n+1}$.
Otrzymujemy funkcję $x=\bigcup_ns_n\in\N^\N$.
Dla każdego $n$ zachodzi $x\in[s_n]\subseteq D_n$,
a zarazem $x\in[s]$. Nie użyliśmy tu dowolnego wyboru
nieskończenie wielu świadków: w każdym kroku wskazaliśmy
pierwszego świadka z ustalonego przeliczalnego porządku.

Gdyby $[s]$ był zawarty w $\bigcup_nF_n$ dla zbiorów
nigdziegęstych, zbiory $D_n=X\setminus\overline F_n$
byłyby otwarte i gęste. Skonstruowany punkt leżałby
w $[s]$, ale poza każdym $F_n$, co daje sprzeczność.
\end{proof}

\subsection{Gra Banacha--Mazura}
Zapis $s\sqsubseteq t$ oznacza, że skończony ciąg $s$
jest początkowym fragmentem ciągu $t$.
\begin{definicja}[Gra Banacha--Mazura nad $A$]
\label{def:ad-banach-mazur}
Niech $A\subseteq\N^\N$. Gracze I i II na przemian
wybierają niepuste skończone ciągi $s_0,s_1,\ldots$
tak, że $s_n$ jest właściwym początkowym fragmentem
$s_{n+1}$. Powstaje jednoznacznie $x=\bigcup_ns_n\in\N^\N$.
Gracz I wygrywa wtedy i tylko wtedy, gdy $x\in A$.
Grę tę oznaczamy przez $\operatorname{BM}(A)$.
\end{definicja}
Każdy ruch jest zakodowany liczbą naturalną, ponieważ zbiór
skończonych ciągów liczb naturalnych jest przeliczalny.
Nielegalny ruch oznacza przegraną gracza, który go wykonał.
W konsekwencji z $\mathrm{AD}$ wynika determinacja tej gry:
dopuszczalne pozycje i reguła przegranej dają grę na
liczbach naturalnych, do której stosuje się definicja AD.

\begin{lemat}[Charakteryzacja strategii przez kategorię]
\label{lem:ad-banach-mazur-category}
Niech $A\subseteq\N^\N$.
\begin{enumerate}
\item Jeżeli gracz II ma strategię wygrywającą w
$\operatorname{BM}(A)$, to $A$ jest pierwszej kategorii.
\item Jeżeli gracz I ma strategię wygrywającą, to istnieje
niepusty cylinder $[s]$ taki, że $[s]\setminus A$
jest pierwszej kategorii w $\N^\N$.
\end{enumerate}
\end{lemat}
\begin{proof}
\emph{Etap 1: strategia gracza II.}
Niech $\tau$ będzie strategią wygrywającą II. Rozpatrzmy
wszystkie skończone legalne historie $h$ zgodne z $\tau$
i kończące się ruchem gracza II; dołączmy też historię pustą.
Jest ich co najwyżej przeliczalnie wiele. Niech $t_h$
będzie ostatnim ciągiem historii $h$ (dla historii pustej
przyjmujemy $[t_h]=\N^\N$).
Niech $V_h$ będzie sumą cylindrów $[u]$, gdzie $u$
jest odpowiedzią $\tau$ na jakiś dozwolony ruch gracza I
zaczynający się po $h$. Wówczas $V_h$ jest zbiorem
otwartym i gęstym w $[t_h]$: w dowolnym mniejszym
cylindrze $[v]\subseteq[t_h]$ gracz I może zagrać
odpowiednio długie przedłużenie $v$, a $\tau$ odpowie
jeszcze dłuższym ciągiem $u\supseteq v$.
Zatem $D_h=(\N^\N\setminus[t_h])\cup V_h$ jest
otwarty i gęsty w $\N^\N$.

Niech $D=\bigcap_hD_h$. Dopełnienie $D$ jest pierwszej
kategorii, gdyż historie indeksowane są przeliczalnie.
Dla $x\in D$ odtwarzamy zgodną z $\tau$ rozgrywkę,
w której wszystkie wybrane ciągi są początkowymi fragmentami
$x$: przynależność $x\in V_h$ pozwala wybrać odpowiedni
ruch I, a $\tau$ daje odpowiedź; pierwszą taką możliwość
bierzemy z ustalonej enumeracji skończonych ciągów.
Wynikiem gry jest $x$. Ponieważ $\tau$ wygrywa,
$x\notin A$, czyli $A\subseteq\N^\N\setminus D$.
Stąd $A$ jest pierwszej kategorii.

\emph{Etap 2: strategia gracza I.}
Niech $\sigma$ będzie strategią wygrywającą gracza I
i niech $s=\sigma(\varnothing)$ będzie jej pierwszym ruchem.
Ograniczmy uwagę do cylindra $[s]$. Po każdym skończonym
fragmencie gry zgodnym z $\sigma$ i zakończonym ruchem I
gracz II może wykonać dowolnie długie przedłużenie,
a $\sigma$ odpowie jeszcze dłuższym ciągiem.
Powtarzając konstrukcję otwartych i gęstych zbiorów
z etapu 1, lecz rozpoczynając od $[s]$ i zamieniając
role graczy, otrzymujemy zbiór rezydualny $D\subseteq[s]$
taki, że każdy punkt $x\in D$ jest wynikiem pewnej
rozgrywki zgodnej z $\sigma$. Wszystkie te punkty leżą
w $A$, bo $\sigma$ jest wygrywająca. Zatem $[s]\setminus A$
jest pierwszej kategorii w $[s]$, a ponieważ $[s]$ jest
otwarto-domknięty, także w całym $\N^\N$.
\end{proof}

\begin{twierdzenie}[AD implikuje własność Baire'a]
\label{thm:ad-baire-property}
W teorii $\mathrm{ZF}+\mathrm{AD}$ każdy podzbiór
$\N^\N$ ma własność Baire'a. To samo zachodzi dla
każdego podzbioru $\R$ w jej zwykłej topologii.
\end{twierdzenie}
\begin{proof}
\emph{Krok 1: część otwarta dla przestrzeni Baire'a.}
Ustalmy $A\subseteq\N^\N$. Niech
\[
U=\bigcup\{[s]:[s]\setminus A\text{ jest pierwszej kategorii}\}.
\]
Jest to zbiór otwarty. Ponieważ baza cylindrów jest
przeliczalna, $U\setminus A$ jest pierwszej kategorii:
stanowi podzbiór przeliczalnej sumy zbiorów $[s]\setminus A$.

\emph{Krok 2: część poza $U$.}
Zagrajmy w $\operatorname{BM}(A\setminus U)$.
Gdyby gracz I miał strategię wygrywającą, Lemat
\ref{lem:ad-banach-mazur-category} dałby niepusty
cylinder $[s]$ taki, że $[s]\setminus(A\setminus U)$
jest pierwszej kategorii. Ponieważ
$[s]\setminus A\subseteq[s]\setminus(A\setminus U)$,
mielibyśmy $[s]\subseteq U$. Wtedy
$[s]\setminus(A\setminus U)=[s]$ byłby pierwszej
kategorii, wbrew Lematowi~\ref{lem:ad-baire-category}.
Zatem I nie ma strategii wygrywającej. Z $\mathrm{AD}$
wynika, że wygrywa II, a Lemat
\ref{lem:ad-banach-mazur-category} pokazuje, że
$A\setminus U$ jest pierwszej kategorii.
Obie strony różnicy symetrycznej $A\triangle U$ są zatem
pierwszej kategorii.

\emph{Krok 3: przejście do prostej.}
Przestrzeń niewymiernych liczb rzeczywistych
$\R\setminus\Q$ jest homeomorficzna z $\N^\N$.
Aby zrozumieć tę identyfikację, najpierw ograniczmy się
do liczb niewymiernych z $(0,1)$. Każda taka liczba
ma jednoznaczne nieskończone rozwinięcie na ułamek
łańcuchowy $[0;a_0,a_1,\ldots]$, gdzie $a_n\geq1$.
Ucięcie rozwinięcia po skończenie wielu współczynnikach
wyznacza przedział możliwych dalszych wartości;
przedziały te tworzą bazę topologii niewymiernych
w $(0,1)$. Odwzorowanie przypisujące liczbie ciąg
$(a_0-1,a_1-1,\ldots)$ i odwzorowanie odwrotne
zachowują więc zbieżność ciągów: dostatecznie długi
wspólny początek współczynników oznacza małą różnicę
liczb, a dla liczby niewymiernej dostatecznie mała
zmiana wartości nie zmienia ustalonego początkowego
fragmentu jej rozwinięcia. Są zatem ciągłe w obie strony.
Całą $\R\setminus\Q$ rozbijamy na przeliczalnie wiele
otwarto-domkniętych podprzestrzeni $(n,n+1)\setminus\Q$,
$n\in\Z$, i stosujemy powyższą konstrukcję w każdej
z nich. Ponieważ $\N\times\N^\N$ jest homeomorficzne
z $\N^\N$, otrzymujemy żądany homeomorfizm.
Homeomorfizm zachowuje zbiory pierwszej kategorii,
więc każda część $A\cap(\R\setminus\Q)$ ma własność
Baire'a w przestrzeni liczb niewymiernych. Ponieważ
$\R\setminus\Q$ jest gęstym zbiorem $G_\delta$
w $\R$, różnica zbioru $A$ od odpowiedniego zbioru
otwartego w $\R$ jest pierwszej kategorii po dołączeniu
zbioru $\Q$. Otrzymujemy żądany wynik dla $\R$.
\end{proof}


\begin{twierdzenie}[AD i własność zbioru doskonałego]
\label{thm:ad-perfect-set}
W teorii ZF+$\mathrm{AD}$ każdy nieprzeliczalny
podzbiór przestrzeni Cantora $2^\omega$,
przestrzeni Baire'a $\omega^\omega$
lub prostej rzeczywistej $\mathbb R$
zawiera niepusty zbiór doskonały
(czyli domknięty, bez punktów izolowanych).
\end{twierdzenie}
\begin{proof}
\emph{Krok 1: gra o podzbiór doskonały.}
Najpierw ustalmy $A\subseteq2^\omega$.
Rozważmy grę, w której gracz I
w swoim $n$-tym ruchu podaje
dowolny skończony ciąg binarny $s_n$,
a gracz II odpowiada jednym bitem
$b_n\in\{0,1\}$.
Wynik rozgrywki jest ciągiem
\[
x=s_0^\frown(b_0)^\frown s_1^\frown(b_1)
          ^\frown s_2^\frown(b_2)^\frown\cdots.
\]
Jest to nieskończony ciąg, ponieważ
każdy ruch gracza II dodaje jeden bit.
Gracz I wygrywa wtedy i tylko wtedy,
gdy $x\in A$.
Zbiór skończonych ciągów binarnych
można jawnie ponumerować liczbami naturalnymi;
nielegalne numery lub ruchy
gracza II spoza $\{0,1\}$ oznaczają
przegraną gracza, który je wykonał.
Wobec tego gra jest kodowana jako gra
o ruchach w $\omega$ i podlega AD.

\emph{Krok 2: strategia II wymusza przeliczalność.}
Załóżmy, że $\tau$ jest wygrywającą strategią
gracza II.
Niech $H$ będzie zbiorem wszystkich
skończonych legalnych historii
zgodnych z $\tau$, kończących się
ruchem gracza II, wraz z historią pustą.
Dla $h\in H$ niech $p_h$
oznacza wypisany dotąd ciąg binarny,
a $\tau(h,s)\in\{0,1\}$ odpowiedź II
na kolejny ruch $s$ gracza I.
Zbiór $H$ jest przeliczalny,
gdyż należy do zbioru skończonych
ciągów liczb naturalnych.

Dla $h\in H$ połóżmy
\[
E_h:=\{x\in2^\omega:
 p_h\sqsubseteq x\ \text{oraz}\
 \forall s\in2^{<\omega}\quad
 p_h^\frown s^\frown(\tau(h,s))
              \not\sqsubseteq x\}.
\]
Każdy $E_h$ ma \emph{co najwyżej jeden}
element. Rzeczywiście, jeśli $x\in E_h$,
to następny bit po $p_h$ jest równy
$1-\tau(h,\varnothing)$.
Gdy znamy pierwszych $k$ bitów $u$
ciągu $x$ po $p_h$, jego następny bit
musi być równy $1-\tau(h,u)$.
Indukcja po $k$ wyznacza więc
wszystkie bity $x$ jednoznacznie.

Pokażemy, że
$A\subseteq\bigcup_{h\in H}E_h$.
Gdyby dla pewnego $x\in A$
nie zachodziło $x\in E_h$ dla żadnego
$h$, to zaczynając od historii pustej
moglibyśmy rekurencyjnie
wybrać pierwszy w ustalonej enumeracji
ruch $s_n$ gracza I, dla którego
\[
p_h^\frown s_n^\frown(\tau(h,s_n))
\sqsubseteq x.
\]
Taki ruch istnieje, bo $x$ przedłuża
bieżące $p_h$ i nie należy do $E_h$.
Odpowiedź $\tau$ przedłuża prefiks
o co najmniej jeden bit;
otrzymalibyśmy zatem pełną rozgrywkę
zgodną z $\tau$ o wyniku $x\in A$.
Przeczy to wygrywającemu
charakterowi $\tau$.
Zatem wskazane zawieranie zachodzi.
Ponieważ $H$ jest przeliczalny
i każdy $E_h$ zawiera najwyżej
jeden element, zbiór $A$ jest
co najwyżej przeliczalny:
przypiszmy każdemu $x\in A$
pierwszą w ustalonej enumeracji
historię $h$ z $x\in E_h$.
Jest to iniekcja do $\omega$,
bez użycia AC.

\emph{Krok 3: strategia I daje zbiór doskonały.}
Jeżeli $A$ nie jest przeliczalny,
gracz II nie ma strategii wygrywającej.
Z AD gracz I ma zatem
strategię wygrywającą $\sigma$.
Dla $b=(b_n)_{n\in\omega}\in2^\omega$
niech $f(b)$ będzie wynikiem gry,
w której II podaje kolejne bity $b_n$,
a I gra zgodnie ze strategią $\sigma$.
Funkcja $f:2^\omega\to2^\omega$
jest ciągła: aby poznać skończony
prefiks jej wartości, wystarcza
poznać skończenie wiele bitów
gracza II, ponieważ każdy z nich
wydłuża wynik o co najmniej jeden bit.
Jest też iniekcją: jeśli
$b\ne b'$, to w pierwszej
współrzędnej, w której się różnią,
poprzednie historie I są jednakowe,
ale do wyniku zostają wpisane
dwa różne bity.

Przestrzeń $2^\omega$ jest zwarta
już w ZF. Istotnie, gdyby pokrycie
otwarte nie miało skończonego
podpokrycia, drzewo cylindrów,
których nie można skończenie pokryć,
byłoby nieskończone. W każdym węźle
co najmniej jedno z dwóch dzieci
również miałoby tę własność.
Wybierając zawsze lewe takie
dziecko, gdy istnieje,
otrzymalibyśmy ciąg $b\in2^\omega$,
którego żadne podstawowe otoczenie
nie można skończenie pokryć,
sprzecznie z tym, że $b$ należy
do pewnego otwartego elementu pokrycia.

Obraz $P:=f[2^\omega]$ jest więc
zwarty, a w przestrzeni Hausdorffa
$2^\omega$ także domknięty.
Nie ma punktów izolowanych:
dla $f(b)\in P$ i jego
dowolnego otoczenia $O$ ciągłość $f$
daje cylinder $[u]\ni b$
odwzorowywany w $O$.
Istnieje $b'\ne b$ o tym samym
prefiksie $u$, a iniektywność
daje $f(b')\ne f(b)$.
Zatem $P$ jest niepusty i doskonały.
Ponieważ $\sigma$ jest wygrywająca,
$P\subseteq A$.

\emph{Krok 4: przeniesienie do
przestrzeni Baire'a i na prostą.}
Definiujemy homeomorfizm
$\theta:\omega^\omega\to D\subseteq2^\omega$,
gdzie $D$ jest zbiorem ciągów binarnych
z nieskończenie wieloma jedynkami,
przez kodowanie
\[
\theta(n_0,n_1,\ldots)
=0^{n_0}1\,0^{n_1}1\,0^{n_2}1\cdots.
\]
Z długości kolejnych bloków zer
można odzyskać wszystkie $n_k$;
obie funkcje są ciągłe,
ponieważ skończony fragment wyniku
zależy od skończenie wielu cyfr wejścia.
Dla nieprzeliczalnego
$B\subseteq\omega^\omega$
zbiór $\theta[B]\subseteq2^\omega$
jest nieprzeliczalny. Z kroków 1--3
zawiera niepusty doskonały,
a w szczególności zwarty zbiór $P$.
Zbiór $\theta^{-1}[P]$
jest zwarty i bez punktów
izolowanych, gdyż $\theta$
jest homeomorfizmem na $D$.
Jako zwarty podzbiór przestrzeni
Hausdorffa jest również domknięty,
a więc doskonały.

Na koniec zauważmy, że
$\mathbb R\setminus\mathbb Q$
jest homeomorficzny z $\omega^\omega$,
co opisaliśmy przy ułamkach łańcuchowych
w dowodzie Twierdzenia~\ref{thm:ad-baire-property}.
Jeśli $C\subseteq\mathbb R$ jest
nieprzeliczalny, to
$C\setminus\mathbb Q$ także
jest nieprzeliczalny:
w przeciwnym razie $C$ byłby
sumą dwóch przeliczalnych zbiorów,
co w ZF nie wymaga nieskończonego
wyboru.
Przenosząc przez homeomorfizm
otrzymany zwarty zbiór doskonały,
uzyskujemy niepusty doskonały
podzbiór $C$.
\end{proof}

\begin{przyklad}[Zbiór Vitaliego nie ma własności Baire'a]
\label{ex:ad-vitali}
Przypomnijmy relację $x\sim y$ wtedy i tylko wtedy, gdy
$x-y\in\Q$. Gdyby zbiór $V$ zawierał dokładnie jednego
reprezentanta każdej klasy, to przesunięcia $V+q$
($q\in\Q$) byłyby parami rozłączne i ich suma dawałaby
całą prostą. Jeśli $V$ byłby pierwszej kategorii, także
$\R=\bigcup_{q\in\Q}(V+q)$ byłaby pierwszej kategorii,
co przeczy twierdzeniu Baire'a. Jeśli $V$ miałby własność
Baire'a i nie był pierwszej kategorii, byłby rezydualny
na pewnym niepustym przedziale $I$. Dla dostatecznie
małego niezerowego $q\in\Q$ przedziały $I$ i $I-q$
przecinają się. W ich przecięciu zbiory $V$ i $V+q$
byłyby jednocześnie rezydualne, więc miałyby wspólny
punkt, sprzecznie z jednoznacznością reprezentantów.
Dlatego w $\mathrm{ZF}+\mathrm{AD}$ taki $V$ nie istnieje.
Planowana dalsza część skryptu poświęcona
ilorazowi $\R/\Q$ omawia silniejsze
ograniczenia wyboru. Nie traktujemy tu
takiej konsekwencji jako udowodnionej
bez osobnego argumentu.
\end{przyklad}

\subsection{Determinacja projektywna a pełne AD}
\begin{definicja}[Zbiory projektywne i PD]
\label{def:ad-projective}
Zbiorami borelowskimi w $\R^n$ są elementy najmniejszej
$\sigma$-algebry zawierającej zbiory otwarte.
Zbiór $A\subseteq\R^n$ jest \emph{analityczny}, jeśli
jest obrazem zbioru borelowskiego w $\R^{n+1}$ przez
rzut na pierwszych $n$ współrzędnych; po przekształceniach
standardowych równoważnie można brać rzuty zbiorów
domkniętych w $\R^n\times\N^\N$.
\emph{Zbiory projektywne} otrzymujemy ze zbiorów
borelowskich przez skończenie wiele operacji dopełnienia
i rzutowania, a więc również przez skończenie wiele
kwantyfikacji po liczbach rzeczywistych.
Aksjomat \emph{determinacji projektywnej} $\mathrm{PD}$
głosi, że każda gra $G(A)$, której zbiór wypłat
$A\subseteq\N^\N$ jest projektywny, jest zdeterminowana.
\end{definicja}
\begin{uwaga}[Granica udowodnionych rezultatów]
\label{rem:ad-deep-regularity}
Klasyczne twierdzenia teorii determinacji głoszą, że
w $\mathrm{ZF}+\mathrm{AD}$ każdy zbiór liczb rzeczywistych
jest mierzalny w sensie Lebesgue'a i ma własność zbioru
doskonałego: jest przeliczalny albo zawiera niepusty
zbiór doskonały. Pod $\mathrm{PD}$ analogiczne własności
przysługują wszystkim \emph{projektywnym} podzbiorom prostej rzeczywistej.
Dowód własności Baire'a przy AD podaliśmy powyżej.
Dowód własności zbioru doskonałego przy AD
podaliśmy w Twierdzeniu~\ref{thm:ad-perfect-set}
w bieżącym rozdziale.
Wyprowadzenie mierzalności z AD znajdziemy dalej w
Twierdzeniu~\ref{thm:ad-measurability}; używa ono
jawnie wydzielonego twierdzenia o uniwersalnej
mierzalności zbiorów analitycznych. Dowodu tego
ostatniego twierdzenia nie powielamy.
Nie wolno też przenosić następstw AD na wszystkie
podzbiory zbioru liczb rzeczywistych przy założeniach
$\mathrm{ZFC}+\mathrm{PD}$:
PD obejmuje tylko zbiory projektywne.
\end{uwaga}


\section{Determinacja a niemożliwość dobrego uporządkowania zbioru liczb rzeczywistych}
\label{sec:ad-no-wellorder-reals}
\begin{twierdzenie}[ZF+AD]
\label{thm:ad-no-wellorder-reals}
Zakładając ZF i aksjomat determinacji AD, zbiór liczb
rzeczywistych $\R$ nie może być dobrze uporządkowany.
\end{twierdzenie}
\begin{proof}
Z wcześniejszego twierdzenia o własności Baire'a pod AD
wynika, że każdy podzbiór $\R$ ma własność Baire'a.
Przypuśćmy, że istnieje dobre uporządkowanie $\R$.
Wówczas istnieje funkcja wybierająca po jednym elemencie
z każdej klasy relacji $x\sim y\iff x-y\in\Q$:
wybieramy element najmniejszy w danym dobrym porządku.
Niech $V\subseteq\R$ będzie zbiorem wybranych reprezentantów.
Pokażemy, że $V$ nie ma własności Baire'a.

Jeśli $V$ jest zbiorem pierwszej kategorii, to również
każdy jego wymierny przekład $V+q$ jest pierwszej kategorii.
Ponieważ $\R=\bigcup_{q\in\Q}(V+q)$, otrzymujemy,
że $\R$ jest zbiorem pierwszej kategorii, wbrew twierdzeniu
Baire'a dla zupełnej przestrzeni metrycznej $\R$.

Jeśli $V$ nie jest zbiorem pierwszej kategorii i ma własność
Baire'a, to istnieje niepusty przedział otwarty $I$,
w którym $V$ jest współpierwszej kategorii: $I\setminus V$
jest pierwszej kategorii. Wybierzmy niezerowe $q\in\Q$
tak małe, aby $I\cap(I+q)$ było niepuste. Zarówno
$V$, jak i $V+q$ są współpierwszej kategorii w tym
przecięciu, więc ich przecięcie jest niepuste, gdyż
niepusty przedział nie jest zbiorem pierwszej kategorii.
Jest to sprzeczne z konstrukcją $V$: dwa różne elementy
jednej klasy modulo $\Q$ nie mogą należeć do $V$.
Obie możliwości prowadzą do sprzeczności.
\end{proof}
\begin{uwaga}
Twierdzenie nie mówi, że nie można dobrze uporządkować
żadnego podzbioru $\R$. Na przykład każdy skończony
podzbiór $\R$ ma dobre uporządkowanie w ZF.
Wynik dotyczy całego zbioru liczb rzeczywistych.
\end{uwaga}


\section{Mierzalność w sensie Lebesgue'a z aksjomatu determinacji}
\label{sec:ad-lebesgue-measurability}
Wcześniej udowodniliśmy, że aksjomat determinacji
zapewnia własność Baire'a wszystkich podzbiorów
zbioru liczb rzeczywistych. To nie jest jednak
synonim mierzalności w sensie Lebesgue'a: zbiór
pierwszej kategorii może mieć dodatnią miarę,
a zbiór miary zero może być gęsty.
Poniżej przeprowadzimy niezależny argument
z użyciem \emph{gry pokryciowej}. Konstrukcję miary Lebesgue\'a,
regularność miary i kryterium Carathéodory\'ego przywołujemy
z odrębnego skryptu \emph{Teoria miary}; istotnym nowym
argumentem w tym rozdziale jest zastosowanie determinacji
do odpowiedniej gry.

\begin{definicja}[Miara zewnętrzna Lebesgue'a i mierzalność]
\label{def:ad-lebesgue}
Dla $A\subseteq\R$ przez $\lambda^*(A)$ oznaczamy
infimum liczb $\sum_{n\in\omega}(b_n-a_n)$
po wszystkich przeliczalnych pokryciach
$A\subseteq\bigcup_{n\in\omega}(a_n,b_n)$
przedziałami otwartymi o skończonych końcach.
Zbiór $E\subseteq\R$ nazywamy \emph{mierzalnym
w sensie Lebesgue'a}, jeżeli dla każdego
$T\subseteq\R$ zachodzi warunek Carathéodory'ego
\[
 \lambda^*(T)=\lambda^*(T\cap E)+\lambda^*(T\setminus E).
\]
Miara Lebesgue'a $\lambda(E)$ zbioru mierzalnego
jest równa $\lambda^*(E)$.
Zbiór $E$ nazywamy \emph{zbiorem miary zero},
jeżeli $\lambda^*(E)=0$. Definicja nie wymaga
uprzedniego założenia, że $E$ jest mierzalny:
każdy zbiór miary zewnętrznej zero jest mierzalny.
\end{definicja}

\begin{uwaga}[Dokładne odsyłacze do skryptu \emph{Teoria miary}]
\label{rem:ad-measure-theory-links}
W niniejszym rozdziale nie rekonstruujemy od podstaw
konstrukcji miary Lebesgue'a ani jej teorii regularności.
Dowody odpowiednich faktów znajdują się w skrypcie
\emph{Teoria miary} z Biblioteki matematycznej, w
następujących miejscach:
\begin{itemize}
\item rozdział 3, \emph{Miary zewnętrzne i konstrukcja
Carathéodory'ego}: kryterium mierzalności Carathéodory'ego
(etykieta źródłowa \nolinkurl{thm:caratheodory});
\item rozdział 4, \emph{Miara Lebesgue'a na przestrzeni
euklidesowej}: następujące wyniki, ze wskazaniem etykiet źródłowych:
\begin{itemize}
\item istnienie miary — \nolinkurl{thm:lebesgue-exists};
\item regularność zewnętrzna — \nolinkurl{thm:lebesgue-outer-regular};
\item regularność wewnętrzna — \nolinkurl{thm:lebesgue-inner-regular};
\item zupełność i aproksymacja borelowska — \nolinkurl{thm:lebesgue-completion}.
\end{itemize}
\end{itemize}
Odsyłacze identyfikują rozdziały i etykiety źródłowe,
lecz nie są poleceniami \texttt{\string\ref}:
kompilacja tego skryptu nie wymaga plików pomocniczych
z odrębnej książki. W szczególności regularność
zewnętrzna jest przywoływanym wynikiem teorii miary,
a nie twierdzeniem udowadnianym ponownie w tym rozdziale.
\end{uwaga}

\begin{przyklad}[Zbiór wymiernych liczb jest miary zero]
Zbiór $\Q$ jest przeliczalny. Ustalmy enumerację
$\Q=\{q_n:n\in\omega\}$. Dla $\varepsilon>0$
przedział
$J_n=(q_n-\varepsilon 2^{-n-3},
q_n+\varepsilon 2^{-n-3})$
ma długość $\varepsilon 2^{-n-2}$.
Suma długości wszystkich $J_n$ wynosi
$\varepsilon/2$, a przedziały pokrywają $\Q$.
Dlatego $\lambda^*(\Q)\le\varepsilon/2$
dla każdego $\varepsilon>0$, więc
$\lambda^*(\Q)=0$. Mimo to $\Q$ jest gęsty w $\R$.
Gęstość i dodatnia miara są więc różnymi własnościami.
\end{przyklad}

\begin{lemat}[Kodowanie pokryć zbioru miary zero]
\label{lem:ad-null-cover-coding}
Pracujemy w $\mathrm{ZF}+\mathrm{DC}$.
Niech $Z\subseteq[0,1]$ będzie zbiorem miary zero.
Dla $\varepsilon>0$ i $n\in\omega$ niech
$\mathcal K_n$ oznacza zbiór wszystkich skończonych
sum przedziałów otwartych o końcach wymiernych,
których łączna długość jest mniejsza niż
$\varepsilon\,4^{-(n+1)}$ (dopuszczamy sumę pustą).
Wtedy istnieją zbiory $H_n\in\mathcal K_n$ takie, że
$Z\subseteq\bigcup_{n\in\omega}H_n$.
Każdy $\mathcal K_n$ jest przeliczalny i posiada
ustaloną enumerację bez potrzeby stosowania
aksjomatu wyboru do nieprzeliczalnej rodziny.
\end{lemat}
\begin{proof}
Z definicji miary zewnętrznej zero możemy wybrać
przeliczalną rodzinę przedziałów otwartych
$(I_j)_{j\in\omega}$ o wymiernych końcach,
pokrywającą $Z$ i spełniającą
$\sum_j |I_j|<\varepsilon/6$.
W tym miejscu $|I_j|$ oznacza \emph{długość przedziału},
a nie moc zbioru; użycie tego samego symbolu
jest lokalnie objaśnione. Powiększenie końców do
liczb wymiernych jest możliwe przy zachowaniu
nierówności, gdyż $\Q$ jest gęsty w $\R$.

Podzielimy ciąg przedziałów na kolejne skończone
bloki. Na początku suma długości wszystkich
przedziałów jest mniejsza od $\varepsilon/4$.
Wybierzmy $N_0$ tak duże, żeby suma długości
ogona o indeksach $j>N_0$ była mniejsza niż
$\varepsilon/16$. Niech $H_0$ będzie sumą
$I_0,\dots,I_{N_0}$. Ponieważ długość sumy
jest nie większa niż suma długości jej składników,
mamy $H_0\in\mathcal K_0$.

Przypuśćmy, że wybraliśmy $N_{n-1}$ tak, aby
suma długości przedziałów $I_j$ z indeksami
$j>N_{n-1}$ była mniejsza niż
$\varepsilon 4^{-(n+1)}$.
Wybieramy $N_n>N_{n-1}$ tak duże, aby
suma długości ogona po $N_n$ była mniejsza niż
$\varepsilon 4^{-(n+2)}$.
Zbiór $H_n=\bigcup_{N_{n-1}<j\le N_n}I_j$
należy do $\mathcal K_n$. Indeksy $N_n$ mogą być
wybierane jako \emph{najmniejsze} indeksy
spełniające wymaganą nierówność, toteż cały
proces jest zdefiniowany przez rekurencję na
$\omega$; nie wymaga dodatkowej funkcji wyboru.
Ponieważ $N_n\to\infty$, wszystkie przedziały
$I_j$ należą do pewnego $H_n$, więc zbiory
$H_n$ pokrywają $Z$.
\end{proof}

\begin{uwaga}[Jeden klasyczny wynik z deskryptywnej teorii mnogości]
\label{rem:ad-analytic-input}
W poniższym dowodzie stosujemy \emph{twierdzenie
o uniwersalnej mierzalności zbiorów analitycznych}:
każdy ciągły obraz przestrzeni Baire'a $\omega^\omega$
w prostej rzeczywistej jest mierzalny w sensie
Lebesgue'a. Jest to samodzielne, nietrywialne
twierdzenie deskryptywnej teorii mnogości,
którego pełen dowód przez operację Suslina
wymaga osobnego opracowania. Nie uznajemy go
za udowodnione w niniejszym skrypcie, ale nie
przypisujemy też jego dowodu książce \emph{Teoria miary}:
w aktualnej wersji tej książki są dowody regularności
miary Lebesgue'a, natomiast nie ma osobnego dowodu
twierdzenia o uniwersalnej mierzalności zbiorów analitycznych.
W literaturze jest to również łączone z twierdzeniem
Łuzina o zbiorach analitycznych; nie należy mylić go
z twierdzeniem Łuzina o aproksymacji funkcji mierzalnych.
Jako szczegółowy tekst źródłowy można wykorzystać
D.~L. Cohna, \emph{Measure Theory}, rozdział 8,
część dotyczącą przestrzeni polskich i zbiorów analitycznych.
To zewnętrzne twierdzenie jest wyraźnie oddzielone od
argumentów teorii mnogości prowadzonych poniżej.
Następny dowód jest \emph{kompletnym wyprowadzeniem
następstwa AD z tego wyraźnie wskazanego twierdzenia}.
\end{uwaga}

\begin{lemat}[Lemat gry pokryciowej]
\label{lem:ad-covering-game}
Pracujemy w $\mathrm{ZF}+\mathrm{DC}+\mathrm{AD}$
i przyjmujemy twierdzenie o mierzalności zbiorów
analitycznych z Uwagi~\ref{rem:ad-analytic-input}.
Niech $S\subseteq[0,1]$ będzie zbiorem takim, że
każdy mierzalny w sensie Lebesgue'a podzbiór
$Z\subseteq S$ ma miarę zero. Wtedy
$\lambda^*(S)=0$.
\end{lemat}
\begin{proof}
\emph{Etap 1: ustalamy grę.}
Niech $\varepsilon>0$. Dla każdego $n$ ustalmy
enumerację $\mathcal K_n=\{K_{n,k}:k\in\omega\}$
z Lematu~\ref{lem:ad-null-cover-coding}.
Gracz I na swoim $n$-tym ruchu wybiera cyfrę
$a_n\in\{0,1\}$, a gracz II wybiera
liczbę naturalną $b_n$. Z cyfr gracza I
powstaje liczba rzeczywista
\[
 x(a)=\sum_{n\in\omega}\frac{a_n}{2^{n+1}}
 \in[0,1].
\]
Gracz I wygrywa dokładnie wtedy, gdy
$x(a)\in S$ oraz
$x(a)\notin\bigcup_{n\in\omega}K_{n,b_n}$.
W przeciwnym razie wygrywa gracz II.
Jest to gra z ruchami w $\omega$ i jasno
określonym zbiorem wypłat, więc AD zapewnia
strategię wygrywającą jednemu z graczy.

\emph{Etap 2: gracz I nie ma strategii wygrywającej.}
Przypuśćmy przeciwnie, że $\sigma$ jest taką
strategią. Po ustaleniu ciągu wszystkich ruchów
$b=(b_0,b_1,\ldots)\in\omega^\omega$ gracza II,
strategia $\sigma$ wyznacza jednoznacznie
wszystkie cyfry gracza I, a więc liczbę
$f_\sigma(b)=x(a)$.
Funkcja $f_\sigma:\omega^\omega\to[0,1]$
jest ciągła: pierwszych $m$ cyfr jej wartości
zależy tylko od pierwszych $m$ ruchów przeciwnika,
a przybliżenie liczby rzeczywistej przez pierwszych
$m$ cyfr binarnych ma błąd najwyżej $2^{-m}$.
Zatem jej obraz $Z=f_\sigma[\omega^\omega]$
jest zbiorem analitycznym, a więc mierzalnym
na mocy wskazanego twierdzenia o mierzalności zbiorów analitycznych.
Ponieważ $\sigma$ wygrywa przeciw każdemu $b$,
każda wartość $f_\sigma(b)$ należy do $S$.
Dlatego $Z\subseteq S$ i z hipotezy
$\lambda(Z)=0$.

Lemat~\ref{lem:ad-null-cover-coding} dostarcza
ciągu $(H_n)$ z $H_n\in\mathcal K_n$
pokrywającego cały $Z$. Gracz II może więc
ustalić ruchy $b_n$ tak, aby $K_{n,b_n}=H_n$.
Strategia $\sigma$ w odpowiedzi wyprodukuje
liczbę rzeczywistą $f_\sigma(b)\in Z$, która
jest pokryta przez pewien $H_n$. Taka rozgrywka
jest przegrana przez gracza I, sprzecznie z
wygrywającym charakterem $\sigma$.

\emph{Etap 3: gracz II wygrywa i pokrywa $S$.}
Z AD wynika, że gracz II ma strategię wygrywającą
$\tau$. Dla każdego skończonego ciągu binarnego
$s=(a_0,\ldots,a_n)$ oznaczmy przez $H_s$
zbiór $K_{n,b_n}$ wybrany przez $\tau$ po ruchach
$(a_0,\ldots,a_n)$ gracza I i odpowiadających
im wcześniejszych ruchach gracza II.
Wówczas
\[
 \lambda(H_s)<\varepsilon 4^{-(n+1)}.
\]
Weźmy dowolną liczbę rzeczywistą $x\in S$.
Wybierzmy jej rozwinięcie binarne $a$,
przy liczbach dyadycznych mniejszych od $1$ wybierając
rozwinięcie kończące się zerami, a dla $x=1$ wybierając
ciąg samych jedynek. Jeśli gracz I podaje
kolejne cyfry tego rozwinięcia, gracz II
postępujący według $\tau$ musi wygrać.
Ponieważ $x\in S$, oznacza to, że
$x\in H_{a\upharpoonright(n+1)}$ dla pewnego $n$.
W konsekwencji
\[
 S\subseteq\bigcup_{n\in\omega}\ \bigcup_{s\in\{0,1\}^{n+1}}H_s.
\]

\emph{Etap 4: obliczamy miarę pokrycia.}
Dla ustalonego $n$ istnieje dokładnie
$2^{n+1}$ ciągów binarnych długości $n+1$.
Podaddytywność miary zewnętrznej daje
\[
\lambda^*(S)
\le\sum_{n\in\omega}2^{n+1}\varepsilon 4^{-(n+1)}
=\varepsilon\sum_{n\in\omega}2^{-(n+1)}
=\varepsilon.
\]
Ponieważ $\varepsilon>0$ było dowolne,
$\lambda^*(S)=0$.
\end{proof}

\begin{twierdzenie}[AD implikuje mierzalność zbiorów liczb rzeczywistych]
\label{thm:ad-measurability}
W teorii $\mathrm{ZF}+\mathrm{DC}+\mathrm{AD}$,
korzystając z klasycznego twierdzenia o mierzalności
zbiorów analitycznych, dla każdego
$A\subseteq\R$ otrzymujemy, że $A$ jest
mierzalny w sensie Lebesgue'a.
\end{twierdzenie}
\begin{proof}
\emph{Etap 1: ograniczamy się do $[0,1]$.}
Niech $X\subseteq[0,1]$. Z regularności zewnętrznej miary Lebesgue\'a
(zob. Uwaga~\ref{rem:ad-measure-theory-links})
istnieją zbiory otwarte
$U_n\supseteq X$ o miarach
$\lambda(U_n)\le\lambda^*(X)+1/(n+1)$.
Zasada wyboru zależnego pozwala ustalić taki
ciąg pokryć. Niech $H=\bigcap_n U_n$.
Jest to zbiór borelowski zawierający $X$.
Monotoniczność miary zewnętrznej i nierówność
$\lambda(H)\le\lambda(U_n)$ dają
\[
 \lambda(H)=\lambda^*(X).
\]

\emph{Etap 2: dopełnienie w obwiedni ma miarę zero.}
Niech $S=H\setminus X$. Jeżeli $Z\subseteq S$
jest mierzalny, to $H\setminus Z$ jest mierzalnym
zbiorem zawierającym $X$. Gdyby $\lambda(Z)>0$,
to ponieważ $\lambda(H)\le1$, otrzymalibyśmy
\[
\lambda(H\setminus Z)=\lambda(H)-\lambda(Z)
<\lambda(H)=\lambda^*(X),
\]
sprzecznie z definicją miary zewnętrznej.
Zatem każdy mierzalny podzbiór $S$ ma miarę zero.
Z Lematu~\ref{lem:ad-covering-game} wynika
$\lambda^*(S)=0$.

\emph{Etap 3: wyprowadzamy mierzalność.}
Zbiory miary zewnętrznej zero są mierzalne,
a miara Lebesgue'a jest zupełna
(zob. Uwaga~\ref{rem:ad-measure-theory-links}). Dlatego
$X=H\setminus S$ jest mierzalny.
Dla dowolnego $A\subseteq\R$ stosujemy
argument do przekształconych liniowo zbiorów
$A\cap[n,n+1]$, $n\in\Z$. Każdy taki fragment
jest mierzalny, więc ich przeliczalna suma,
równa $A$, również jest mierzalna.
\end{proof}

\section{Co naprawdę wynika z determinacji projektywnej?}
\label{sec:pd-regularity}
\begin{twierdzenie}[Regularność zbiorów projektywnych z PD]
\label{thm:pd-regularity}
Przy założeniu $\mathrm{ZFC}+\mathrm{PD}$ każdy
projektywny podzbiór zbioru liczb rzeczywistych
jest mierzalny w sensie Lebesgue'a, ma własność
Baire'a i jest przeliczalny albo zawiera
niepusty zbiór doskonały.
W dowodzie mierzalności korzystamy ze wskazanego
klasycznego twierdzenia o zbiorach analitycznych.
\end{twierdzenie}
\begin{proof}
W dowodach mierzalności
(Twierdzenie~\ref{thm:ad-measurability}),
własności Baire'a
(Twierdzenie~\ref{thm:ad-baire-property})
i własności zbioru doskonałego
(Twierdzenie~\ref{thm:ad-perfect-set})
aksjomat AD był stosowany do gier zbudowanych
ze zbioru wypłat $A$ przez ustalone kodowania,
operacje borelowskie i przeciwobrazy przez
funkcje ciągłe. Jeśli $A$ jest projektywny,
tak otrzymane zbiory wypłat również są projektywne:
klasy projektywne są zamknięte na dopełnienie,
rzutowanie, operacje borelowskie i ciągłe
przeciwobrazy. Zatem do każdej z tych gier
wystarcza już PD. Pozostałe części dowodów
— konstrukcja pokryć miary zero, wyprowadzenie
sprzeczności ze strategią i budowanie
zbioru doskonałego — nie wymagają determinacji
innych gier. Otrzymujemy trzy własności
regularności dla $A$.
\end{proof}

\begin{uwaga}[Nie mylić PD z pełnym AD]
Z $\mathrm{ZFC}+\mathrm{PD}$ nie wynika, że
\emph{każdy} podzbiór zbioru liczb rzeczywistych
ma własność Baire'a lub jest mierzalny.
Pełny aksjomat wyboru nadal pozwala skonstruować
zbiory Vitaliego; ich istnienie nie przeczy PD,
ponieważ nie są zbiorami projektywnymi.
W szczególności nie wolno zastąpić w tezie
poprzedniego twierdzenia określenia
„każdy zbiór projektywny” przez
„każdy podzbiór zbioru liczb rzeczywistych”.
\end{uwaga}


\section{Determinacja projektywna a duże liczby kardynalne}
\label{sec:ad-pd-large-cardinal-strength}
W tej sekcji odróżniamy trzy rodzaje twierdzeń:
wyniki wynikające bezpośrednio z podanych
wcześniej definicji, głębokie twierdzenia o
modelach wewnętrznych oraz konsekwencje takich
twierdzeń, które możemy dowieść na podstawie
wcześniejszych rozdziałów. Pracujemy w otaczającej
teorii ZFC, chyba że wyraźnie wskazano inną teorię.

\begin{definicja}[Model $L(\mathbb R)$ i determinacja wewnętrzna]
\label{def:lr-ad}
Niech $V$ będzie otaczającym uniwersum zbiorów,
a $\R^V$ jego zbiorem liczb rzeczywistych.
Przez $L(\R)$ oznaczamy najmniejszy model
wewnętrzny ZF zawierający wszystkie liczby
porządkowe i wszystkie liczby rzeczywiste
z $V$, konstruowany przez iterację operacji
\emph{definiowalnych podzbiorów} z parametrem $\R^V$.
Zdanie $L(\R)\models\mathrm{AD}$ oznacza, że
\emph{każda gra o długości $\omega$ na liczbach
naturalnych, której zbiór wypłat jest elementem
$L(\R)$, jest zdeterminowana w tym modelu}.
Nie oznacza to, że $V\models\mathrm{AD}$.
\end{definicja}

\begin{uwaga}[Nie utożsamiać zbioru liczb rzeczywistych z jego kodowaniami]
Dla celów semantyki gry używamy przestrzeni
$\omega^\omega$ wszystkich funkcji $\omega\to\omega$.
Jej elementy są ciągami liczb naturalnych;
zbiór $\R$ ma inną definicję. W ZFC przestrzenie
te mają wzajemnie jednoznaczne kodowania borelowskie,
co pozwala przenosić klasy projektywne oraz
własności determinacji w ich standardowym ujęciu.
\end{uwaga}

\begin{definicja}[Liczba kardynalna Woodina — sformułowanie]
\label{def:woodin-cardinal}
Niech $\delta$ będzie nieprzeliczalną liczbą
kardynalną. Mówimy, że $\delta$ jest
\emph{liczbą kardynalną Woodina}, jeżeli dla
każdej funkcji $f:\delta\to\delta$ istnieje
$\kappa<\delta$ i elementarne zanurzenie
$j:V\to M$ do tranzytywnego modelu klasowego $M$
takie, że $\kappa$ jest punktem krytycznym $j$,
$f[\kappa]\subseteq\kappa$ oraz
$V_{j(f)(\kappa)}\subseteq M$.
Punktem krytycznym elementarnego zanurzenia jest
najmniejsza liczba porządkowa $\alpha$ spełniająca
$j(\alpha)\ne\alpha$. Zapis $f[\kappa]$ oznacza
obraz zbioru $\kappa$ przy funkcji $f$.
Sformułowanie używa semantyki klas i wymaga
precyzyjnej interpretacji w stosownej teorii klas;
jego równoważne warianty za pomocą extenderów
należą do teorii dużych liczb kardynalnych.
\end{definicja}

\begin{uwaga}[Granica samodzielności definicji]
Samo podanie definicji liczby kardynalnej Woodina
nie stanowi dowodu istnienia takiej liczby
kardynalnej ani nie daje determinacji projektywnej.
Nie twierdzimy też, że wystarczy istnienie jednej
liczby kardynalnej Woodina do udowodnienia pełnej PD.
\end{uwaga}

\begin{twierdzenie}[Martin--Steel: determinacja projektywna]
\label{thm:martin-steel-projective}
W teorii ZFC z założenia istnienia nieskończenie
wielu liczb kardynalnych Woodina wynika
aksjomat determinacji projektywnej $\mathrm{PD}$.
Jest to głębokie twierdzenie teorii modeli
wewnętrznych, przywołane tutaj ze wskazaniem
źródła, a nie wyprowadzone z definicji.
\end{twierdzenie}
\begin{uwaga}[Źródło głębokiego twierdzenia]
Opis zależności, ograniczeń i literatury znajduje się
w Stanford Encyclopedia of Philosophy,
\emph{Large Cardinals and Determinacy}:
\url{https://plato.stanford.edu/entries/large-cardinals-determinacy/}.
W książce rozwijamy szczegółowo podstawową teorię
mierzalnych liczb kardynalnych, natomiast pełny
dowód powyższego twierdzenia wymaga maszynerii
iterowalnych modeli wewnętrznych oraz extenderów.
\end{uwaga}

\begin{twierdzenie}[Martin--Steel--Woodin: determinacja w $L(\R)$]
\label{thm:martin-steel-woodin-lr}
Pracujemy w ZFC. Jeżeli istnieje rosnący ciąg
$(\delta_n)_{n\in\omega}$ liczb kardynalnych
Woodina oraz mierzalna liczba kardynalna
$\kappa$ spełniająca $\delta_n<\kappa$ dla
każdego $n\in\omega$, to
\[
 L(\R)\models\mathrm{AD}.
\]
Jest to przywołane twierdzenie głębokie;
nie zastępujemy jego dowodu stwierdzeniem,
że wynika \emph{bezpośrednio} z istnienia
mierzalnej liczby kardynalnej.
\end{twierdzenie}
\begin{uwaga}[Źródło i znaczenie hipotez]
Zob. również rozdział 33 podręcznika Jecha
\emph{Set Theory}, \emph{Determinacy},
\url{https://fa.ewi.tudelft.nl/~hart/onderwijs/set_theory/Jech/33-determinacy.pdf}.
Hipoteza mówi, że \emph{jedna} mierzalna liczba
kardynalna leży powyżej całej wskazanej rodziny
liczb kardynalnych Woodina. Nie wynika to z samej
informacji, że istnieje nieskończenie wiele
liczb kardynalnych Woodina.
\end{uwaga}

\begin{lemat}[Prawdziwość formuł projektywnych w modelach o tych samych liczbach rzeczywistych]
\label{lem:projective-same-reals}
Niech $M\subseteq V$ będzie tranzytywnym modelem
wewnętrznym ZF o tych samych liczbach naturalnych
i rzeczywistych co $V$, czyli
$\omega^M=\omega^V$ i $\R^M=\R^V$.
Niech $A\subseteq\omega^\omega$ będzie zbiorem
projektywnym zadanym formułą z parametrami będącymi
liczbami rzeczywistymi. Wtedy dla każdego
$a\in\omega^\omega$ formuła definiująca
przynależność $a\in A$ ma tę samą wartość logiczną
w $M$ i w $V$.
\end{lemat}
\begin{proof}
Zbiory borelowskie definiowane za pomocą kodów
borelowskich złożonych z liczb naturalnych są
absolutne między modelami o tych samych liczbach
rzeczywistych: relacja przynależności do zbioru
kodowanego jest obliczana przez dobrze ufundowaną
rekurencję po drzewie kodu, które należy do obu
modeli, a zbiory liczb naturalnych są te same.
Prowadzimy następnie indukcję po liczbie operacji
projektywnych. Dopełnienie zmienia prawdę
na fałsz po obu stronach, a projekcja jest
kwantyfikacją po elementach $\omega^\omega$.
Modele mają dokładnie te same takie elementy,
więc kwantyfikator egzystencjalny daje tę samą
wartość logiczną. W ten sposób wszystkie
skończone etapy definicji projektywnej są
absolutne. Nie używamy tutaj aksjomatu wyboru
\emph{wewnątrz} modelu $M$.
\end{proof}

\begin{wniosek}[Determinacja w $L(\R)$ implikuje PD w otaczającym uniwersum]
\label{cor:lr-ad-implies-pd}
Jeżeli $V\models\mathrm{ZFC}$ oraz
$L(\R)\models\mathrm{AD}$, to $V\models\mathrm{PD}$.
\end{wniosek}
\begin{proof}
Niech $A\subseteq\omega^\omega$ będzie zbiorem
projektywnym z parametrami będącymi liczbami
rzeczywistymi z $V$. Ponieważ $L(\R)$ zawiera
wszystkie liczby rzeczywiste i porządkowe,
zawiera także kod definiujący $A$, a dzięki
wyróżnianiu oraz projektywnej definicji zawiera
sam zbiór $A$.
Z założenia $L(\R)\models\mathrm{AD}$
w $L(\R)$ istnieje strategia wygrywająca jednego
z dwóch graczy w grze o wypłacie $A$. Strategia
jest funkcją na zbiorze wszystkich skończonych
ciągów liczb naturalnych; można ją zakodować
jako liczbę rzeczywistą, toteż ta sama strategia
należy także do $V$. Z Lematu~\ref{lem:projective-same-reals}
wynika, że $A$ ma identyczne elementy w obu
modelach. Dlatego każda rozgrywka zgodna ze
strategią kończy się zwycięstwem tego samego
gracza także w $V$. Gra jest zdeterminowana,
a ponieważ $A$ był dowolny, zachodzi PD.
\end{proof}

\begin{wniosek}[Regularność projektywnych podzbiorów liczb rzeczywistych]
\label{cor:woodin-projective-regularity}
Przy hipotezach Twierdzenia~\ref{thm:martin-steel-woodin-lr}
w otaczającym uniwersum ZFC każdy projektywny
podzbiór zbioru liczb rzeczywistych jest
mierzalny w sensie Lebesgue'a, ma własność
Baire'a i jest przeliczalny albo zawiera
niepusty zbiór doskonały.
\end{wniosek}
\begin{proof}
Twierdzenie~\ref{thm:martin-steel-woodin-lr}
daje $L(\R)\models\mathrm{AD}$. Na mocy
Wniosku~\ref{cor:lr-ad-implies-pd} zachodzi
PD w otaczającym uniwersum ZFC.
Twierdzenie~\ref{thm:pd-regularity} zastosowane
do dowolnego projektywnego podzbioru zbioru
liczb rzeczywistych daje wszystkie trzy własności.
\end{proof}

\begin{uwaga}[Dlaczego nie otrzymujemy pełnego AD w ZFC?]
Z aksjomatu wyboru w $V$ wynika istnienie
gier niezdeterminowanych
(Twierdzenie~\ref{thm:ac-not-ad}).
Dlatego w powyższym twierdzeniu celowo
piszemy $L(\R)\models\mathrm{AD}$, a nie
$V\models\mathrm{AD}$. To pierwsze zdanie jest
zgodne z aksjomatem wyboru w otaczającym
uniwersum, drugie — nie.
\end{uwaga}


\section{Redukcje Wadge'a i dalsza konsekwencja determinacji projektywnej}
\label{sec:pd-wadge-reductions}
W tej sekcji pracujemy w ZFC, a dodatkowo przyjmujemy
aksjomat determinacji projektywnej PD tylko tam, gdzie
został wyraźnie wymieniony. Nie używamy tu twierdzeń
o mierzalności zbiorów liczb rzeczywistych.

\begin{definicja}[Redukowalność Wadge'a]
\label{def:wadge-reducibility}
Niech $A,B$ będą podzbiorami przestrzeni Baire'a
$\omega^\omega$, czyli zbioru wszystkich ciągów liczb
naturalnych z topologią generowaną przez cylindry
$[s]$ z Definicji~\ref{def:baire-cylinders}.
Piszemy $A\leq_{\mathrm W}B$, jeżeli istnieje
funkcja ciągła $f:\omega^\omega\to\omega^\omega$
spełniająca
\[
\forall x\in\omega^\omega\qquad
   x\in A\ \Longleftrightarrow\ f(x)\in B.
\]
Oznacza to równość $A=f^{-1}[B]$.
Nie wymagamy, aby $f$ była iniekcją, surjekcją
ani bijekcją. Relacja $\leq_{\mathrm W}$ jest
preporządkiem: jest zwrotna (funkcja identycznościowa)
i przechodnia (złożenie funkcji ciągłych).
\end{definicja}

\begin{lemat}[Zbiory projektywne a kodowanie gier Wadge'a]
\label{lem:wadge-payoff-projective}
Niech $A,B\subseteq\omega^\omega$ będą zbiorami
projektywnymi, w sensie Definicji~\ref{def:ad-projective}.
W grze, w której gracz I podaje po jednej liczbie
naturalnej, a gracz II w każdej turze podaje liczbę
naturalną albo symbol pominięcia ruchu, zbiór
rozgrywek wygrywanych przez gracza II według
następujących reguł jest projektywny:
\begin{enumerate}
\item gracz II musi podać nieskończenie wiele liczb,
      aby utworzyć ciąg $y\in\omega^\omega$;
\item po oznaczeniu przez $x\in\omega^\omega$ ciągu
      liczb gracza I wymagamy równoważności
      $x\in A\iff y\in B$.
\end{enumerate}
\end{lemat}
\begin{proof}
Kodujemy jeden ruch parą liczb naturalnych,
przy czym drugi składnik jest albo symbolem pominięcia,
albo liczbą naturalną. Przy ustalonym bijektywnym
kodowaniu par całe rozgrywki tworzą przestrzeń
$\omega^\omega$. Ciąg $x$ zależy ciągle od rozgrywki.
Warunek, że II podał nieskończenie wiele liczb,
jest warunkiem borelowskim (dla każdego $n$ istnieje
późniejszy ruch niebędący pominięciem). Na zbiorze
rozgrywek spełniających ten warunek ciąg $y$ zależy
ciągle od całej rozgrywki: pierwsze $k$ jego wyrazów
są wyznaczone przez pewien skończony początkowy
fragment rozgrywki, którego długość może zależeć
od rozgrywki i od $k$. Przeciwobrazy zbiorów
projektywnych przez takie kodowania pozostają
projektywne; własność ta wynika przez indukcję
z definicji klas projektywnych, ponieważ przeciwobrazy
zbiorów borelowskich przez funkcje ciągłe są borelowskie,
a rzutowanie i dopełnienie zachowują hierarchię.
Równoważność dwóch zdań zapisujemy przez skończone
sumy, przecięcia i dopełnienia. Po uwzględnieniu
borelowskiego warunku legalności gry otrzymujemy
zbiór projektywny.
\end{proof}

\begin{twierdzenie}[Półliniowość porządku Wadge'a dla zbiorów projektywnych]
\label{thm:pd-projective-wadge}
Zakładając ZFC i PD, dla dowolnych dwóch
projektywnych zbiorów $A,B\subseteq\omega^\omega$
zachodzi co najmniej jedna z możliwości
\[
 A\leq_{\mathrm W}B
 \qquad\text{lub}\qquad
 B\leq_{\mathrm W}(\omega^\omega\setminus A).
\]
W szczególności nie chodzi o zwykłą porównywalność
bez dopełnienia, ani o porównanie mocy zbiorów.
\end{twierdzenie}
\begin{proof}
Rozważmy grę opisaną w
Lemacie~\ref{lem:wadge-payoff-projective}.
W każdej turze I podaje liczbę naturalną, a II
może podać liczbę naturalną albo pominąć ruch.
Jeśli II poda tylko skończenie wiele liczb,
przegrywa. Jeżeli poda nieskończenie wiele liczb,
to wygrywa dokładnie wtedy, gdy uzyskane ciągi
$x$ i $y$ spełniają $x\in A\iff y\in B$.
Na mocy lematu zbiór wypłat jest projektywny,
zatem PD gwarantuje strategię wygrywającą jednego
z graczy.

\emph{Przypadek pierwszy: zwycięską strategię ma II.}
Dla ustalonego wejścia $x\in\omega^\omega$
pozwalamy I kolejno podawać jego wyrazy,
a II odpowiada zgodnie ze swoją strategią.
Strategia II jest wygrywająca, więc w każdej takiej
rozgrywce podaje nieskończenie wiele liczb.
Oznaczmy utworzony ciąg przez $f(x)\in\omega^\omega$.
Aby wykazać ciągłość $f$ w punkcie $x$, ustalmy
liczbę naturalną $k$. Po pewnej skończonej liczbie
ruchów strategia II wytworzyła pierwszych $k$
wyrazów $f(x)$. Każde wejście o tym samym
wystarczająco długim prefiksie wywołuje te same
odpowiedzi strategii do tego momentu. Zatem
pierwszych $k$ wyrazów $f(x)$ jest lokalnie
stałych, co jest kryterium ciągłości w przestrzeni
Baire'a. Z warunku zwycięstwa II mamy
$x\in A\iff f(x)\in B$ dla każdego $x$.
Otrzymaliśmy $A\leq_{\mathrm W}B$.

\emph{Przypadek drugi: zwycięską strategię ma I.}
Pozwalamy graczowi II podawać w każdej turze
jedną liczbę naturalną i nigdy nie pomijać ruchu.
Dla dowolnego ciągu $y\in\omega^\omega$ jego
kolejne wyrazy stanowią taką legalną rozgrywkę
gracza II. Strategia gracza I wyznacza w odpowiedzi
ciąg $g(y)=x\in\omega^\omega$. Ponieważ każdy
skończony prefiks $g(y)$ zależy jedynie od
skończonej liczby wyrazów $y$, funkcja $g$ jest
ciągła. Gracz I wygrywa przeciw każdej rozgrywce
II, w tym przeciw wybranej, dlatego
\[
 \neg\bigl(g(y)\in A\iff y\in B\bigr).
\]
Powyższy warunek jest równoważny
$y\in B\iff g(y)\notin A$. W konsekwencji
$B=g^{-1}[\omega^\omega\setminus A]$,
czyli $B\leq_{\mathrm W}(\omega^\omega\setminus A)$.
Oba przypadki wyczerpują możliwości.
\end{proof}

\begin{przyklad}[Dlaczego potrzebne jest dopełnienie?]
Niech $z\in\omega^\omega$ będzie ciągiem stale
równym zeru, $A=\{z\}$ oraz
$B=\{y\in\omega^\omega:y(0)=0\}$. Zbiór $A$
jest domknięty, ale nieotwarty, natomiast $B$
jest jednocześnie otwarty i domknięty. Gdyby
$A=f^{-1}[B]$ dla funkcji ciągłej $f$, to
$A$ byłby otwarty, ponieważ przeciwobraz zbioru
otwartego przez funkcję ciągłą jest otwarty.
Tak więc $A\not\leq_{\mathrm W}B$.

Mimo to $B\leq_{\mathrm W}(\omega^\omega\setminus A)$:
niech $u\in\omega^\omega$ będzie ciągiem stale
równym jedności i połóżmy $g(y)=u$, gdy
$y(0)=0$, oraz $g(y)=z$ w przeciwnym razie.
Funkcja $g$ jest ciągła, bo ma stałe wartości
na dwóch rozłącznych zbiorach otwarto-domkniętych,
a dokładnie na $B$ przyjmuje wartość spoza $A$.
Ten przykład ilustruje matematyczne znaczenie
dołączenia dopełnienia w tezie twierdzenia.
\end{przyklad}

\begin{uwaga}[Zakres twierdzenia i literatura]
Twierdzenie~\ref{thm:pd-projective-wadge}
ma samodzielny dowód oparty na PD. Głębsze
rezultaty, takie jak determinacja gier o wyższej
długości, struktura rang Wadge'a i twierdzenia
o skalach projektywnych, wymagają oddzielnych
dowodów. Por. D. Marker, \emph{Descriptive Set Theory},
\url{https://www.math.uic.edu/~marker/math512/dst.pdf},
oraz Stanford Encyclopedia of Philosophy,
\emph{Large Cardinals and Determinacy},
\url{https://plato.stanford.edu/entries/large-cardinals-determinacy/}.
\end{uwaga}


\section{Pierwsze stopnie hierarchii Wadge'a i pełne znaczenie gry redukcyjnej}
\label{sec:wadge-first-levels}
Z poprzedniej sekcji wiemy, że determinacja projektywna
PD daje półliniowość redukowalności Wadge'a dla zbiorów
projektywnych (Twierdzenie~\ref{thm:pd-projective-wadge}).
Teraz zbierzemy konstrukcje, które można w całości
udowodnić w ZFC, bez zakładania PD. Odróżnimy trzy
pojęcia: \emph{redukowalność ciągłą}, \emph{równoważność
Wadge'a} i \emph{przynależność do klasy borelowskiej}.

\begin{definicja}[Równoważność Wadge'a i zupełność dla klasy zbiorów]
\label{def:wadge-equivalence-completeness}
Dla dowolnych zbiorów $A,B\subseteq\omega^\omega$
piszemy $A\equiv_{\mathrm W}B$, jeżeli
$A\leq_{\mathrm W}B$ i $B\leq_{\mathrm W}A$.
Klasa równoważności tej relacji jest nazywana
\emph{stopniem Wadge'a}. Jeśli $\Gamma$ jest
rodziną podzbiorów $\omega^\omega$, mówimy, że
$B\subseteq\omega^\omega$ jest \emph{zupełny
względem redukcji Wadge'a dla $\Gamma$}, jeżeli
$B\in\Gamma$ i każdy $A\in\Gamma$ spełnia
$A\leq_{\mathrm W}B$. Nie chodzi tutaj o
zupełność metryczną ani o moc zbioru.
\end{definicja}

\begin{lemat}[Jakie operacje zachowuje redukcja ciągła?]
\label{lem:wadge-continuous-preimages}
Niech $A,B\subseteq\omega^\omega$ i
$A\leq_{\mathrm W}B$. Jeśli $B$ jest otwarty,
to $A$ jest otwarty; jeśli $B$ jest domknięty,
to $A$ jest domknięty. Jeśli $B$ jest przeliczalną
sumą zbiorów domkniętych, to $A$ również ma tę
postać. Analogicznie przeciwobrazy przeliczalnych
przecięć zbiorów otwartych są przeliczalnymi
przecięciami zbiorów otwartych.
\end{lemat}
\begin{proof}
Istnieje funkcja ciągła $f:\omega^\omega\to\omega^\omega$
taka, że $A=f^{-1}[B]$. Z definicji ciągłości
przeciwobraz zbioru otwartego jest otwarty;
przeciwobraz zbioru domkniętego jest domknięty,
ponieważ $f^{-1}[\omega^\omega\setminus B]
=\omega^\omega\setminus f^{-1}[B]$.
Dla dowolnej rodziny $(B_n)_{n\in\omega}$ mamy
\[
 f^{-1}\!\left[\bigcup_{n\in\omega}B_n\right]
 =\bigcup_{n\in\omega}f^{-1}[B_n],
 \quad
 f^{-1}\!\left[\bigcap_{n\in\omega}B_n\right]
 =\bigcap_{n\in\omega}f^{-1}[B_n].
\]
Te równości wynikają bezpośrednio z definicji
przeciwobrazu i uzasadniają obie pozostałe tezy.
\end{proof}

\begin{twierdzenie}[Zbiór otwarty zupełny dla zbiorów otwartych]
\label{thm:wadge-open-complete}
Niech
\[
 U=\{x\in\omega^\omega:\exists n\in\omega\ (x(n)\ne0)\}.
\]
Zbiór $U$ jest otwarty i nie jest domknięty.
Dla dowolnego zbioru otwartego $A\subseteq\omega^\omega$
istnieje funkcja ciągła $f_A:\omega^\omega\to\omega^\omega$
taka, że $A=f_A^{-1}[U]$.
\end{twierdzenie}
\begin{proof}
Otwartość $U$ wynika z równości
$U=\bigcup_{n\in\omega}\{x:x(n)\ne0\}$;
każdy zbiór w tej sumie jest otwarty i domknięty.
Niech $z$ będzie ciągiem stale równym zeru, a
$x^{(m)}$ ciągiem przyjmującym wartość $1$
dokładnie na pozycji $m$ i wartość $0$
na pozostałych pozycjach. Każdy $x^{(m)}\in U$,
a $x^{(m)}\to z\notin U$, więc $U$ nie jest
domknięty.

Ustalmy otwarty zbiór $A$. Dla każdego ciągu
$x\in\omega^\omega$ i każdego $n\in\omega$ połóżmy
\[
 f_A(x)(n)=
 \begin{cases}
  1,&\text{gdy }[x\upharpoonright(n+1)]\subseteq A,\\
  0,&\text{w przeciwnym przypadku}.
 \end{cases}
\]
Symbol $x\upharpoonright(n+1)$ oznacza skończony
ciąg pierwszych $n+1$ wyrazów $x$, a cylindry
$[s]$ zdefiniowaliśmy w Definicji~\ref{def:baire-cylinders}.
Dla ustalonego $n$ wartość $f_A(x)(n)$ zależy
wyłącznie od tego skończonego ciągu, więc każda
współrzędna funkcji $f_A$ jest ciągła. Oznacza to
ciągłość $f_A$ w topologii produktowej.
Jeśli $x\in A$, otwartość daje pewien prefiks
$s=x\upharpoonright k$ z $[s]\subseteq A$.
Dla wszystkich $n+1\ge k$ mamy
$[x\upharpoonright(n+1)]\subseteq[s]\subseteq A$,
więc $f_A(x)\in U$.
Jeśli $x\notin A$, żaden cylinder zawierający
$x$ nie może być podzbiorem $A$; wówczas
$f_A(x)$ jest ciągiem stale równym zeru,
a więc $f_A(x)\notin U$. Dowiedliśmy
$A=f_A^{-1}[U]$.
\end{proof}

\begin{twierdzenie}[Zbiór domknięty zupełny dla zbiorów domkniętych]
\label{thm:wadge-closed-complete}
Niech $K=\{z\}$, gdzie $z\in\omega^\omega$
jest ciągiem stale równym zeru.
Zbiór $K$ jest domknięty i nie jest otwarty.
Dla każdego zbioru domkniętego $F\subseteq\omega^\omega$
istnieje ciągła funkcja $g_F:\omega^\omega\to\omega^\omega$
spełniająca $F=g_F^{-1}[K]$.
\end{twierdzenie}
\begin{proof}
Zbiór $K$ jest domknięty jako przecięcie
$\bigcap_{n\in\omega}\{x:x(n)=0\}$.
Nie jest otwarty, ponieważ każdy cylinder
zawierający $z$ zawiera także ciągi różne od $z$.
Dla ustalonego zbioru domkniętego $F$ definiujemy
\[
 g_F(x)(n)=
 \begin{cases}
  0,&\text{jeśli }[x\upharpoonright(n+1)]\cap F\ne\varnothing,\\
  1,&\text{jeśli }[x\upharpoonright(n+1)]\cap F=\varnothing.
 \end{cases}
\]
Także tutaj wartość każdej współrzędnej zależy
od skończonego początkowego fragmentu $x$,
zatem $g_F$ jest ciągła. Dla $x\in F$
każdy cylinder zawierający $x$ przecina $F$,
więc $g_F(x)=z\in K$.
Jeśli $x\notin F$, otwartość dopełnienia $F$
gwarantuje istnienie cylindra
$[x\upharpoonright k]$ rozłącznego z $F$.
Wtedy $g_F(x)(n)=1$ dla każdego $n+1\ge k$,
a więc $g_F(x)\notin K$.
Otrzymujemy $F=g_F^{-1}[K]$.
\end{proof}

\begin{wniosek}[Pierwsze rozróżnialne stopnie Wadge'a]
\label{cor:wadge-first-degrees}
W ZFC zachodzą następujące fakty:
\begin{enumerate}
\item wszystkie niepuste, właściwe zbiory \emph{jednocześnie otwarte i domknięte} w $\omega^\omega$ są wzajemnie równoważne Wadge'a;
\item jeśli $U$ i $K$ są zbiorami z poprzednich dwóch twierdzeń, to ani $U\leq_{\mathrm W}K$, ani $K\leq_{\mathrm W}U$;
\item mimo to $U\equiv_{\mathrm W}(\omega^\omega\setminus K)$.
\end{enumerate}
\end{wniosek}
\begin{proof}
Dla pierwszego punktu ustalmy niepuste właściwe zbiory jednocześnie otwarte i domknięte $A,B$.
Wybierzmy $u\in B$ oraz $v\notin B$ i określmy $f(x)=u$ dla $x\in A$, zaś $f(x)=v$ dla $x\notin A$.
Dla dowolnego otwartego $O\subseteq\omega^\omega$ przeciwobraz $f^{-1}[O]$ jest jednym ze zbiorów $\varnothing,A,\omega^\omega\setminus A,\omega^\omega$; każdy z nich jest otwarty. Zatem $f$ jest ciągła i $f^{-1}[B]=A$.
Zamiana ról $A,B$ daje równoważność.

Zbiór $U$ jest otwarty, ale nie jest domknięty; zbiór $K$ jest domknięty, ale nie jest otwarty (co wykazano w poprzednich twierdzeniach).
Gdyby $U\leq_{\mathrm W}K$, ciągłość przeciwobrazu wymuszałaby domkniętość $U$; analogicznie $K\leq_{\mathrm W}U$ wymuszałoby otwartość $K$.
Wreszcie $\omega^\omega\setminus K$ jest niepustym właściwym zbiorem otwartym. Z zupełności $U$ dla zbiorów otwartych wynika $\omega^\omega\setminus K\leq_{\mathrm W}U$.
Przeciwną redukcję daje zupełność $\omega^\omega\setminus K$: jest to zbiór otwarty zupełny, ponieważ $K$ jest domknięty zupełny, a dopełnianie zachowuje równanie $f^{-1}[B]=A$.
\end{proof}

\begin{twierdzenie}[Strategie gracza II i funkcje ciągłe]
\label{thm:wadge-strategy-continuity}
Niech $A,B\subseteq\omega^\omega$. W grze Wadge'a
z Lematu~\ref{lem:wadge-payoff-projective},
w której gracz I podaje kolejne wyrazy ciągu $x$,
a gracz II podaje kolejne wyrazy ciągu $y$
lub pomija ruch, przy czym musi ostatecznie podać
nieskończenie wiele wyrazów, zachodzi równoważność:
\[
 A\leq_{\mathrm W}B
 \quad\Longleftrightarrow\quad
 \text{gracz II ma strategię wygrywającą w tej grze}.
\]
Ten wynik nie wymaga założenia determinacji
wszystkich gier.
\end{twierdzenie}
\begin{proof}
\emph{Strategia daje funkcję ciągłą.}
To przejście zostało wykazane w pierwszym przypadku
dowodu Twierdzenia~\ref{thm:pd-projective-wadge}:
dla każdego wejścia $x$ rozgrywamy strategię
gracza II przeciwko kolejnym wyrazom $x$,
a ciąg wyrazów podanych przez II jest wartością
ciągłej funkcji $f(x)$. Warunek zwycięstwa
zapewnia $A=f^{-1}[B]$. Dowód ciągłości
nie wykorzystuje PD, lecz tylko określoną
strategię wygrywającą.

\emph{Funkcja ciągła daje strategię.}
Załóżmy, że $A=f^{-1}[B]$ dla pewnej funkcji
ciągłej $f:\omega^\omega\to\omega^\omega$.
Po każdej turze gracza I znamy jego dotychczasowy
skończony ciąg ruchów $s\in\omega^{<\omega}$.
Możemy wówczas określić największą liczbę
naturalną $m(s)\leq|s|$ taką, że wszystkie
ciągi $f(x)$ przy $x\in[s]$ mają jednakowe
pierwsze $m(s)$ wyrazów. Taka liczba istnieje,
bo $m=0$ zawsze spełnia ten warunek, a rozpatrujemy
skończony zbiór możliwych wartości $m$.

Gracz II przechowuje skończony ciąg $t$ już
podanych wyrazów. Jeśli $|t|<m(s)$, podaje
następny wyraz wspólnego prefiksu długości $m(s)$;
jeśli $|t|\ge m(s)$, pomija ruch.
W każdym kroku jego poprzednie odpowiedzi są
prefiksem $f(x)$ dla każdego $x\in[s]$.
Ustalmy dowolną kompletną rozgrywkę gracza I,
czyli ciąg $x\in\omega^\omega$.
Ciągłość $f$ w punkcie $x$ daje dla każdego
$k\in\omega$ pewną liczbę naturalną $N$ taką,
że wszystkie $x'\in[x\upharpoonright N]$
spełniają $f(x')\upharpoonright k=f(x)\upharpoonright k$.
Dlatego dla wszystkich $n\ge\max(N,k)$
mamy $m(x\upharpoonright(n+1))\ge k$.
Jeśli do tego czasu gracz II nie podał jeszcze
$k$ wyrazów, będzie podawał po jednym wyrazie
w kolejnych turach, aż uzyska długość $k$.
Ponieważ $k$ było dowolne, gracz II podaje
nieskończenie wiele wyrazów, tworząc dokładnie
$f(x)$. Równość $A=f^{-1}[B]$ zapewnia
wygraną gracza II w każdej takiej rozgrywce.
\end{proof}

\begin{uwaga}[Zależności od determinacji projektywnej]
Trzy pierwsze poziomy redukowalności pokazano
konstrukcyjnie w ZFC. Dopiero półliniowość
porządku Wadge'a dla wszystkich zbiorów
projektywnych używa PD
(Twierdzenie~\ref{thm:pd-projective-wadge}).
Nie twierdzimy tutaj, że pełna hierarchia Wadge'a
jest dobrze uporządkowana: jej analiza, rangi
Wadge'a i twierdzenia o skalach wymagają
oddzielnych wyników deskryptywnej teorii mnogości.
\end{uwaga}


\section{Samodualność i pierwsze stopnie redukcji Wadge'a}
\label{sec:wadge-selfduality-clopen}
Zbiór $A\subseteq\omega^\omega$ nazywamy \emph{samodualnym względem redukcji Wadge'a}, jeżeli $A\equiv_{\mathrm W}(\omega^\omega\setminus A)$.
To własność stopnia redukcji, nie równość zbiorów.

\begin{twierdzenie}[Samodualność niebanalnych zbiorów otwarto-domkniętych]
\label{thm:wadge-clopen-selfdual}
Jeśli $A\subseteq\omega^\omega$ jest niepustym, właściwym zbiorem otwartym i domkniętym, to $A\equiv_{\mathrm W}(\omega^\omega\setminus A)$.
Natomiast $\varnothing$ i $\omega^\omega$ nie są samodualne.
\end{twierdzenie}
\begin{proof}
Dopełnienie $A$ również jest niepustym, właściwym zbiorem otwartym i domkniętym. Z pierwszego punktu Wniosku~\ref{cor:wadge-first-degrees} otrzymujemy obie redukcje.
Dla dowolnej funkcji $f:\omega^\omega\to\omega^\omega$ mamy $f^{-1}[\varnothing]=\varnothing$ oraz $f^{-1}[\omega^\omega]=\omega^\omega$. Dlatego $\omega^\omega\not\leq_{\mathrm W}\varnothing$ i $\varnothing\not\leq_{\mathrm W}\omega^\omega$.
\end{proof}

\begin{twierdzenie}[Niesamodualność właściwie otwartego stopnia zupełnego]
\label{thm:wadge-open-nonselfdual}
Jeśli $U\subseteq\omega^\omega$ jest zbiorem otwartym zupełnym z Twierdzenia~\ref{thm:wadge-open-complete}, to $U\not\equiv_{\mathrm W}(\omega^\omega\setminus U)$.
\end{twierdzenie}
\begin{proof}
Dopełnienie $U$ jest domknięte. Gdyby $U\leq_{\mathrm W}(\omega^\omega\setminus U)$, ciągły przeciwobraz zbioru domkniętego byłby domknięty. To przeczy temu, że $U$ nie jest domknięty.
\end{proof}

\begin{przyklad}[Dwa różne typy stopni]
Zbiór $A=\{x\in\omega^\omega:x(0)=0\}$ jest niepusty, właściwy, otwarty i domknięty, a więc samodualny. Zbiór $U=\{x:\exists n\ x(n)=0\}$ jest otwarty, ale nie domknięty: ciągi o coraz późniejszym pierwszym zerze zbiegają do ciągu stale równego $1$, który nie należy do $U$. Ponieważ $U$ jest zupełny dla zbiorów otwartych, jego stopień nie jest samodualny.
\end{przyklad}


\begin{uwaga}[Audyt warunku samodualności]
\label{rem:wadge-selfdual-audit-32}
Samodualność dotyczy równoważności przez \emph{ciągłe przeciwobrazy}, a nie istnienia homeomorfizmu zamieniającego zbiór z dopełnieniem. W szczególności dowód Twierdzenia~\ref{thm:wadge-clopen-selfdual} wymaga, aby obie części przestrzeni były otwarte: tylko wtedy dwuwartościowa funkcja stała na obu częściach jest ciągła. Dla zbioru otwartego, który nie jest domknięty, taki argument nie działa.
\end{uwaga}